\documentclass[12pt,a4paper,UTF8]{ctexart}
\usepackage[a4paper,top=22mm,bottom=22mm,left=22mm,right=22mm,headheight=15pt]{geometry}
\usepackage{amsmath,amssymb,mathtools,bm}
\usepackage{graphicx}
\usepackage{xcolor}
\usepackage{booktabs,longtable,array}
\usepackage{enumitem}
\usepackage{fancyhdr}
\usepackage{caption}
\usepackage{float}
\usepackage{microtype}
\usepackage{setspace}
\usepackage[hidelinks,unicode]{hyperref}
\usepackage[most]{tcolorbox}

\definecolor{ContestBlue}{HTML}{17365D}
\definecolor{SoftBlue}{HTML}{EEF4FA}
\definecolor{SoftGray}{HTML}{F5F6F7}
\hypersetup{pdftitle={第101套},pdfauthor={物理竞赛教学资料}}

\setstretch{1.24}
\setlength{\parindent}{2em}
\setlength{\parskip}{0.35em}
\setlist{itemsep=0.35em,topsep=0.4em}
\captionsetup{font=small,labelfont=bf,justification=centering}
\setcounter{secnumdepth}{0}

\pagestyle{fancy}
\fancyhf{}
\fancyhead[L]{\small 第101套}
\fancyfoot[C]{\thepage}
\renewcommand{\headrulewidth}{0.4pt}

\ctexset{
  section={format=\Large\bfseries\color{ContestBlue},beforeskip=1.2em,afterskip=0.7em},
  subsection={format=\large\bfseries\color{ContestBlue},beforeskip=1.0em,afterskip=0.5em},
  subsubsection={format=\normalsize\bfseries,beforeskip=0.8em,afterskip=0.35em}
}

\makeatletter
\newcommand{\tightlist}{\setlength{\itemsep}{0pt}\setlength{\parskip}{0pt}}
\newcommand{\maxwidth}{\ifdim\Gin@nat@width>\linewidth\linewidth\else\Gin@nat@width\fi}
\newcommand{\maxheight}{\ifdim\Gin@nat@height>0.72\textheight 0.72\textheight\else\Gin@nat@height\fi}
\newcommand{\pandocbounded}[1]{\begingroup\setkeys{Gin}{width=\maxwidth,height=\maxheight,keepaspectratio}#1\endgroup}
\makeatother

\begin{document}

\begin{titlepage}
  \centering
  \vspace*{58mm}
  {\Huge\bfseries\color{ContestBlue} 第101套\par}
  \vfill
  \vspace*{28mm}
\end{titlepage}


\section{第 1 题}\label{ux7b2c-1-ux9898}

在距竖直墙水平距离为 \(l\)、离地高度为 \(h\) 的 \(A\) 点，以大小为
\(v_0\) 的初速度将球斜向地面抛出。球先在距墙水平距离为 \(x\) 的地面点
\(B\) 与地面发生弹性碰撞，再在墙上 \(C\) 点发生弹性碰撞，最后回到 \(A\)
点。

\begin{enumerate}
\def\labelenumi{\arabic{enumi}.}
\tightlist
\item
  将球视为质点，求 \(x\)。
\item
  将球视为质量为 \(m\)、半径为 \(r\) 的均匀实心球。球从 \(A\) 点以速度
  \(v_0\)
  水平抛出，初始角速度为零；球依次与地面、墙面碰撞，且两次碰撞过程中接触点始终相对滑动，滑动摩擦因数均为
  \(\mu\)。若人站在原处接球，手应从 \(A\) 点向上移动多少？
\end{enumerate}

\begin{figure}[H]
\centering
\includegraphics[width=0.694\linewidth,height=\textheight,keepaspectratio,alt={图 1-1　球依次在地面点 B、墙面点 C 弹性碰撞后返回抛出点所在竖线。}]{../figures/question-01.pdf}
\caption*{图 1-1　球依次在地面点 \(B\)、墙面点 \(C\)
弹性碰撞后返回抛出点所在竖线。}
\end{figure}

\subsection{一、题型定位}\label{ux4e00ux9898ux578bux5b9aux4f4d}

按落地前、地面反弹后和墙面反弹后三段处理抛体运动。加入球的自转后，摩擦冲量方向由接触点滑动速度决定，有限半径也会改变接触位置。

\subsection{二、知识点}\label{ux4e8cux77e5ux8bc6ux70b9}

\begin{itemize}
\tightlist
\item
  抛体运动的分段拼接与弹性反射。
\item
  利用``同一高度对应的两个时刻''消去碰撞点高度。
\item
  滑动碰撞中的法向冲量、切向摩擦冲量及速度突变。
\item
  对未知初速度方向建立代数方程，并依据实际运动顺序筛选物理解。
\end{itemize}

\subsection{三、物理图像}\label{ux4e09ux7269ux7406ux56feux50cf}

把运动看成三段抛物线与两个瞬时反射。质点模型中，地面只使竖直速度反向，墙面只使水平速度反向；实心球模型中，法向冲量负责反弹，切向冲量同时改变质心切向速度和角速度。

\subsection{四、解题入口}\label{ux56dbux89e3ux9898ux5165ux53e3}

先把未知初速度分解为水平分量 \(u\) 和向下竖直分量
\(w\)，再引入落地前竖直速率 \(q\)。用同一高度对应的两个时刻写出 \(l-x\)
与 \(l+x\)，这是消去轨迹细节的入口。

先不要急着追求最终公式。第（1）问只做两件事：把落地前后的竖直运动接起来，把三段水平路程接起来；完全弹性碰撞意味着相应速度分量只反向、不改大小。第（2）问再把``瞬间碰撞''改用冲量描述，并逐次核对摩擦冲量究竟改变哪一个速度分量。整题的主线是

\[\text{分段运动}\longrightarrow\text{碰撞处速度突变}\longrightarrow\text{下一段运动}.\]

把两次反射展开后，回程的水平路程为 \(l+x\)；落地前则是
\(l-x\)。两者分别对应同一高度的两个时间根。

\subsection{五、解析}\label{ux4e94ux89e3ux6790}

图 1-2 左侧把球从 \(A\) 经地面碰撞点 \(B\)、墙面碰撞点 \(C\)
后返回的两段抛体轨迹画在同一坐标中；右侧分别画出地面、墙面碰撞时的法向冲量、滑动摩擦冲量和自转方向。第（1）问使用左图的镜像几何，第（2）问按右图逐次更新平动速度与角速度。

\begin{figure}[H]
\centering
\includegraphics[width=0.9\linewidth,height=\textheight,keepaspectratio,alt={图 1-2　两段抛体轨迹及地面、墙面碰撞时的冲量与自转方向。}]{../figures/solution-01.pdf}
\caption*{图 1-2　两段抛体轨迹及地面、墙面碰撞时的冲量与自转方向。}
\end{figure}

\textbf{（1）}

先看图 1-2
左侧。两次弹性法向碰撞可分别用地面镜像和墙面镜像展开，但两碰撞间的水平距离为
\(x\)、第一段为 \(l-x\)，所以必须保留中间碰撞位置再消元。

取水平方向向右、竖直方向向上。设初速度的水平分量为
\(u>0\)，初始竖直分量为 \(-w\)，其中 \(w\geq 0\)，则

\[u^2+w^2=v_0^2.\]

球从 \(A\) 落到地面前的竖直速率为

\[q=\sqrt{w^2+2gh}.\]

下面把这两个量的来源写清楚。球在 \(A\) 点的竖直速度为
\(-w\)，落地前的竖直速度为 \(-q\)。由竖直方向的匀加速关系

\[(-q)^2-(-w)^2=2(-g)(0-h),\]

得到 \(q^2=w^2+2gh\)。若从抛出到落地的时间为 \(t_1\)，则

\[0=h-wt_1-\frac12gt_1^2.\]

整理为关于 \(t_1\) 的二次方程：

\[gt_1^2+2wt_1-2h=0.\]

它的两个代数根为

\[t_1=\frac{-w\pm\sqrt{w^2+2gh}}{g}
=\frac{-w\pm q}{g}.\]

时间必须为正，而 \(q>w\)，所以只能取

\[t_1=\frac{q-w}{g}.\]

由于 \(AB\) 的水平距离为 \(l-x\)，

\[l-x=\frac{u(q-w)}{g}. \tag{1}\]

在 \(B\) 点发生完全弹性碰撞后，竖直速度由 \(-q\) 变为 \(+q\)。球从 \(B\)
经墙面反射后回到 \(A\) 所在竖线，累计水平路程为 \(x+l\)，故该段总时间为

\[T=\frac{l+x}{u}.\]

球从地面以竖直初速度 \(q\) 上升，再次到达高度 \(h\) 时满足

\[h=qT-\frac12gT^2.\]

将其整理为

\[gT^2-2qT+2h=0,\]

两个时间根为

\[\begin{aligned}
T
&=\frac{2q\pm\sqrt{4q^2-8gh}}{2g}\\
&=\frac{q\pm\sqrt{q^2-2gh}}{g}\\
&=\frac{q\pm w}{g}.
\end{aligned}\]

其中 \((q-w)/g\) 对应球反弹后在上升阶段第一次经过高度 \(h\)；\((q+w)/g\)
对应越过最高点后在下降阶段再次经过高度 \(h\)。本题从 \(B\)
到最终接球点的水平路程为 \(l+x\)，大于从 \(A\) 到 \(B\) 的水平路程
\(l-x\)，而两段水平速率同为 \(u\)，所以
\(T>t_1=(q-w)/g\)，必须取较大的根：

\[T=\frac{q+w}{g}.\]

因此

\[l+x=\frac{u(q+w)}{g}. \tag{2}\]

将式（1）、（2）相加：

\[2l=\frac{u[(q-w)+(q+w)]}{g}=\frac{2uq}{g},\]

因此 \(uq=gl\)。再用式（2）减式（1）：

\[2x=\frac{u[(q+w)-(q-w)]}{g}=\frac{2uw}{g},\]

因此

\[uq=gl,\qquad x=\frac{uw}{g}. \tag{3}\]

令 \(z=u^2\)，由 \(q^2=w^2+2gh=v_0^2-z+2gh\) 得

\[z\left(v_0^2-z+2gh\right)=g^2l^2,\]

即

\[z^2-\left(v_0^2+2gh\right)z+g^2l^2=0.\]

所以

\[z=\frac{v_0^2+2gh\pm\sqrt{\left(v_0^2+2gh\right)^2-4g^2l^2}}{2}. \tag{4}\]

由 \(z=u^2\) 及 \(w^2=v_0^2-u^2=v_0^2-z\)，有

\[x^2=\frac{u^2w^2}{g^2}
=\frac{z(v_0^2-z)}{g^2}.\]

因为 \(x\) 表示距离，应取正根，最终

\[\boxed{x=\frac1g\sqrt{z\left(v_0^2-z\right)}}. \tag{5}\]

式（4）若产生两个代数根，应同时检验
\(0<z\leq v_0^2\)、\(0<x<l\)，并保留满足``先落地、再撞墙、最后回到
\(A\)''这一时间顺序的根。

\textbf{代数化简说明：} 式（4）的判别式要求

\[\left(v_0^2+2gh\right)^2-4g^2l^2\geq0,\]

即

\[v_0^2+2gh\geq2gl.\]

这给出了存在相应轨迹的必要条件。得到 \(z\) 后还要检查
\(w^2=v_0^2-z\geq0\)；判别式非负并不自动保证初速度方向真实存在。

也可直接写出

\[u=\sqrt z,\qquad
w=\sqrt{v_0^2-z},\qquad
q=\frac{gl}{\sqrt z},\]

再用

\[t_1=\frac{q-w}{g},\qquad
t_{BC}=\frac{x}{u},\qquad
t_{CA}=\frac{l}{u}\]

逐段检查三个时间均为正。

\emph{【数学技巧：利用速度平方变量降次】} 若直接对 \(u\)
求解会出现四次方程；令 \(z=u^2\) 后，方程立即降为关于 \(z\) 的二次方程。

\emph{【拓展视角：镜像法】} 将墙面后的运动关于墙面镜像，\(C\to A\)
的反射轨迹可看成球继续飞向墙后虚点
\(A'\)。于是地面反弹后的整段轨迹等效为从 \(B\) 到 \(A'\)
的一次完整抛体运动。该图像能直接说明水平总路程为 \(x+l\)，也能避免在
\(C\) 点重新设置一套运动方程。

\textbf{（2）}

先按原解的几何近似，忽略球半径对碰撞位置的影响；此时球心在地面
\(y=0\)、墙面 \(x=l\)
处反弹。题目给出有限半径，下面在近似式之后补出球心接触几何的修正。

记球自由落地前的竖直速率为

\[s=\sqrt{2gh},\qquad t_0=\sqrt{\frac{2h}{g}}.\]

第一次碰撞的法向冲量大小为 \(2ms\)，故滑动摩擦冲量大小为
\(2\mu ms\)。它与球的水平滑动方向相反，所以球离开地面后的水平速度大小为

\[u=v_0-2\mu s. \tag{6}\]

这里用到的是冲量---动量定理。碰前、碰后竖直速度分别为 \(-s\) 与
\(+s\)，所以地面对球的法向冲量为

\[J_{N1}=m[s-(-s)]=2ms.\]

接触点相对地面向右滑动，摩擦冲量向左，大小为
\(J_{f1}=\mu J_{N1}=2\mu ms\)，故

\[m(u-v_0)=-J_{f1}=-2\mu ms,\]

这就得到式（6）。

式（6）还给出一个必须检查的运动条件：球碰地后仍要向墙运动，因此需要

\[u=v_0-2\mu s>0.\]

若
\(u\leq0\)，球在第一次碰撞后已经停止向墙靠近，题目所描述的``随后碰墙''过程便不会发生。

若需同时追踪球的自转，地面摩擦冲量对球心的角冲量为

\[\Delta L=-2\mu msr.\]

均匀实心球的转动惯量为 \(I=\frac25mr^2\)，故地面碰撞后的角速度为

\[\omega_1=\frac{\Delta L}{I}
=-\frac{5\mu s}{r}.\]

这里负号表示按图示取逆时针为正时，球获得顺时针自转。题目已明确两次碰撞期间接触点始终相对滑动，因此摩擦冲量大小可直接取
\(\mu\) 倍法向冲量；但摩擦方向仍须由接触点相对速度判断。

球的落地点距 \(A\) 所在竖线为 \(v_0t_0\)。令落地点至墙的距离为

\[d=l-v_0t_0,\]

因为从 \(A\) 所在竖线到墙的距离为 \(l\)，而球在落地前已向墙运动了
\(v_0t_0\)。只有 \(d>0\)
时，球才是先碰地面、后碰墙面。地面碰撞后水平速度恒为
\(u\)，所以从地面到墙面的飞行时间为

\[\tau=\frac{d}{u}. \tag{7}\]

球撞墙前的竖直速度为

\[v_y^-=s-g\tau.\]

墙面对球的法向冲量大小为 \(2mu\)，相应的滑动摩擦冲量大小为
\(2\mu m u\)。在题设所对应的持续滑动方向下，墙面摩擦冲量向上，故

\[v_y^+=s-g\tau+2\mu u. \tag{8}\]

这一步同样可逐项写为

\[J_{N2}=\left|m[(-u)-u]\right|=2mu,\]

其中只比较垂直墙面的动量改变量；于是切向摩擦冲量为
\(J_{f2}=2\mu m u\)。它向上时，竖直速度增量为

\[\Delta v_y=\frac{J_{f2}}m=2\mu u,\]

从而得到式（8）。

更严谨地说，撞墙前球最右侧接触点相对墙面的切向速度为

\[v_{\mathrm{rel},y}=v_y^-+\omega_1r.\]

只有当该量为负时，墙面对球的滑动摩擦冲量才向上。由此可将所选摩擦方向写成适用条件

\[s-g\tau-5\mu s<0.\]

若该不等式不成立，式（8）中的 \(2\mu u\) 应改为
\(-2\mu u\)。题设``始终有相对滑动''排除了粘连，却不替代这一步方向判断。

球到达墙面时的高度为

\[y_C=s\tau-\frac12g\tau^2.\]

这是以地面为零高度、以离开地面为计时起点写出的竖直位移。墙面反射后，水平方向速度只改变符号，大小仍为
\(u\)。墙到 \(A\) 所在竖线的水平距离为 \(l\)，所以返回时间为

\[t_r=\frac lu.\]

返回阶段的竖直位移为

\[\Delta y_r=v_y^+t_r-\frac12gt_r^2.\]

因此接球高度为

\[y_f=y_C+v_y^+\frac{l}{u}-\frac12g\left(\frac{l}{u}\right)^2.\]

因此手的上移量为

\[\boxed{\Delta h=s\tau-\frac12g\tau^2+\left(s-g\tau+2\mu u\right)\frac{l}{u}-\frac12g\left(\frac{l}{u}\right)^2-h},\]

其中

\[u=v_0-2\mu\sqrt{2gh},\qquad
\tau=\frac{l-v_0\sqrt{2h/g}}{u}.\]

在上述忽略半径的几何近似中，质量 \(m\) 与半径 \(r\)
在结果中消去。该表达式成立的前提是两次碰撞期间接触点的相对滑动方向确与上述冲量方向一致，且
\(u>0\)、\(\tau>0\)。

冲量更新中的 \(m\) 消去，转动对接触速度的贡献以 \(\omega r\)
出现；但有限半径仍会通过实际落差和到墙距离进入飞行时间。

\textbf{有限半径与滑动方向的完整结果。} 以题目 \(A\)
点为初始球心位置。球心接触地面时高度为 \(r\)，接触墙时横坐标为
\(l-r\)。定义

\[H=h-r,\quad L=l-r,\quad s=\sqrt{2gH},\quad t_0=\sqrt{2H/g},\]
\[u=v_0-2\mu s,\quad d=L-v_0t_0,\quad \tau=d/u,\quad t_r=L/u.\]

墙面碰前的接触点切向速度为 \(w_c=s-g\tau-5\mu s\)。令
\(\sigma=-\operatorname{sgn}w_c\)，即摩擦向上时 \(\sigma=+1\)，向下时
\(\sigma=-1\)。手相对初始高度的位移为

\[\boxed{\Delta h=r+s\tau-\frac12g\tau^2+(s-g\tau+2\mu u\sigma)t_r-\frac12gt_r^2-h.}\]

这里 \(h>r\)、\(l>r\)，且必须有 \(d>0\)。地面碰撞结束时的滑动速度是
\(v_0-7\mu s\)，所以全程保持向右滑动要求
\(v_0>7\mu s\)。墙面碰撞结束时的切向滑动速度是
\(w_c+7\mu u\sigma\)；全程保持同一滑动方向要求
\(|w_c|>7\mu u\)。只检查碰前符号不够。还须确认球在抵墙前及返回接球前没有再次触地。若
\(\Delta h<0\)，手应下移，其大小为
\(-\Delta h\)；题设参数不能自动保证上移。让 \(r\to0\) 且取
\(\sigma=+1\)，就回到前面的近似式。

\newpage

\section{第 2 题}\label{ux7b2c-2-ux9898}

一根长度为 \(l\) 的轻质刚杆两端分别固定质量为 \(2m\) 与 \(3m\)
的小球，组成异形哑铃。该哑铃静置在水平光滑桌面上。另有一质量为 \(80m\)
的小球，以速度 \(v_0\) 沿水平方向运动，并在哑铃质心上方距离为 \(x\)
处与哑铃发生完全弹性碰撞。

\begin{enumerate}
\def\labelenumi{\arabic{enumi}.}
\tightlist
\item
  求碰撞后质量为 \(80m\) 的小球及哑铃两端小球的速度。
\item
  若要使质量为 \(80m\) 的小球与哑铃发生第二次碰撞，求 \(x/l\)
  的取值范围。
\item
  求两次碰撞之间的时间；列出一元方程即可。
\end{enumerate}

\begin{figure}[H]
\centering
\includegraphics[width=0.538\linewidth,height=\textheight,keepaspectratio,alt={图 2-1　入射小球在哑铃质心上方 x 处发生横向碰撞。}]{../figures/question-02.pdf}
\caption*{图 2-1　入射小球在哑铃质心上方 \(x\) 处发生横向碰撞。}
\end{figure}

\subsection{一、题型定位}\label{ux4e00ux9898ux578bux5b9aux4f4d-1}

偏心碰撞先用冲量、动量与弹性接触条件求碰后状态。随后转到相对运动，检查小球是否再次碰到有限长杆，并区分杆内接触和端点接触。

\subsection{二、知识点}\label{ux4e8cux77e5ux8bc6ux70b9-1}

\begin{itemize}
\tightlist
\item
  刚体质心位置与转动惯量。
\item
  偏心完全弹性碰撞中的动量、角动量和机械能守恒。
\item
  刚体平面运动中端点速度的合成。
\item
  二次碰撞的追及条件与转动几何约束。
\end{itemize}

\subsection{三、物理图像}\label{ux4e09ux7269ux7406ux56feux50cf-1}

碰撞瞬间，哑铃获得质心平动和绕质心转动。碰后在质心参考系中，小球沿水平直线运动，哑铃则像一根不断转动的指针；第二次碰撞就是这条直线与有限杆段再次相交。

\subsection{四、解题入口}\label{ux56dbux89e3ux9898ux5165ux53e3-1}

先求质心位置和转动惯量，再把碰后状态压缩为三个未知量
\(v\)、\(V_c\)、\(\omega\)。依次写线动量、关于质心的角动量和机械能守恒。

这道题分成两个层次。第一层只研究``碰撞瞬间''：把哑铃看成一个刚体，用动量、关于质心的角动量和机械能三个守恒关系求出碰后状态。第二层才研究``碰撞以后''：在哑铃质心系中，看一条转动的有限杆是否再次扫到小球的直线轨迹。建议学生在纸边始终写着

\[\text{质心平动}+\text{绕质心转动}=\text{刚体上任一点的运动}.\]

先彻底做懂第（1）问，再进入第二次碰撞；否则学生容易把碰撞守恒与碰后追及混成一套方程。

\subsection{五、解析}\label{ux4e94ux89e3ux6790-1}

图 2-2 在哑铃质心系中画出碰后哑铃绕质心转过
\(\omega t\)，同时入射小球沿相对质心的直线运动。第（1）、（2）问先求哑铃质心速度与角速度，第（3）问再令转动后的杆段与小球相对轨迹重新相交，确定第二次碰撞时间。

\begin{figure}[H]
\centering
\includegraphics[width=0.524\linewidth,height=\textheight,keepaspectratio,alt={图 2-2　转动杆段与小球相对质心的直线轨迹再次相交。}]{../figures/solution-02.pdf}
\caption*{图 2-2　转动杆段与小球相对质心的直线轨迹再次相交。}
\end{figure}

\textbf{（1）}

图 2-2 中 \(3m\) 端距质心较近、\(2m\)
端距质心较远，先由质心定义确定两段长度；入射小球的作用线在质心上方
\(x\)，这同时决定碰撞角动量方程的力臂。

哑铃质心到 \(3m\) 端、\(2m\) 端的距离分别为

\[a=\frac{2l}{5},\qquad b=\frac{3l}{5}.\]

若质心到 \(3m\) 端的距离为 \(a\)，则到 \(2m\) 端的距离为
\(l-a\)。由质心定义

\[3m\,a=2m(l-a),\]

解得
\(a=2l/5\)、\(l-a=3l/5\)。较重的一端离质心更近，这是检查距离是否写反的直接方法。

哑铃绕质心的转动惯量为

\[I=3m\left(\frac{2l}{5}\right)^2+2m\left(\frac{3l}{5}\right)^2
=\frac65ml^2. \tag{1}\]

设碰后入射小球速度为 \(v\)，哑铃质心平动速度为 \(V_c\)，角速度为
\(\omega\)。沿碰撞方向动量守恒、对哑铃质心角动量守恒以及总动能守恒分别给出

\[\begin{aligned}
80mv_0&=80mv+5mV_c,\\
80mv_0x&=80mvx+I\omega,\\
\frac12(80m)v_0^2&=\frac12(80m)v^2+\frac12(5m)V_c^2+\frac12I\omega^2.
\end{aligned}\]

第一式把哑铃整体质量写成
\(2m+3m=5m\)。第二式以哑铃质心为矩心：入射小球的作用线到质心的垂直距离为
\(x\)，所以小球碰撞前后的角动量分别为 \(80mv_0x\) 与
\(80mvx\)；哑铃获得的自转角动量为
\(I\omega\)。第三式可用是因为题设给出完全弹性碰撞。

下面按``先消去两个线性未知量，再处理能量方程''的顺序联立。由动量守恒式

\[80v_0=80v+5V_c\]

移项得到

\[V_c=16(v_0-v), \tag{2}\]

以及

\[\omega=\frac{80mx}{I}(v_0-v)
=\frac{200x}{3l^2}(v_0-v). \tag{3}\]

将式（2）、（3）代入能量守恒式，并约去公共因子 \(\frac12m\)：

\[80(v_0^2-v^2)
=5\cdot16^2(v_0-v)^2
+\frac65l^2\left[\frac{200x}{3l^2}(v_0-v)\right]^2.\]

提取 \(v_0-v\)：

\[80(v_0-v)(v_0+v)
=(v_0-v)^2\left(1280+\frac{16000x^2}{3l^2}\right).\]

非平凡碰撞满足 \(v\neq v_0\)，约去 \(v_0-v\)，再通分整理：

\[v_0+v
=\left(16+\frac{200x^2}{3l^2}\right)(v_0-v).\]

为把最后一步的代数也展开，记

\[B=16+\frac{200x^2}{3l^2}.\]

方程 \(v_0+v=B(v_0-v)\) 展开为

\[v_0+v=Bv_0-Bv,\]

把含 \(v\) 的项移到左边、含 \(v_0\) 的项移到右边：

\[(1+B)v=(B-1)v_0.\]

因此

\[v=\frac{B-1}{B+1}v_0
=\frac{45l^2+200x^2}{51l^2+200x^2}v_0.\]

为简化后续书写，定义

\[D=51l^2+200x^2.\]

于是

\[\boxed{v=\frac{45l^2+200x^2}{D}v_0}.\]

把它代回式（2）：

\[\begin{aligned}
V_c
&=16(v_0-v)\\
&=16v_0\left(1-\frac{45l^2+200x^2}{D}\right)\\
&=16v_0\frac{6l^2}{D}
=\boxed{\frac{96l^2}{D}v_0}.
\end{aligned}\]

再代回式（3）：

\[\begin{aligned}
\omega
&=\frac{200x}{3l^2}(v_0-v)\\
&=\frac{200x}{3l^2}v_0\frac{6l^2}{D}\\
&=\boxed{\frac{400x}{D}v_0}.
\end{aligned}\]

最后求两个端点的速度。碰撞瞬间杆竖直，绕质心转动产生的端点速度沿水平方向，且两端方向相反。\(2m\)
端距质心 \(3l/5\)，故

\[\begin{aligned}
V_{2m}
&=V_c-\frac{3l}{5}\omega\\
&=\frac{96l^2}{D}v_0
-\frac{3l}{5}\frac{400x}{D}v_0\\
&=\boxed{\frac{96l^2-240lx}{D}v_0}.
\end{aligned}\]

\(3m\) 端距质心 \(2l/5\)，故

\[\begin{aligned}
V_{3m}
&=V_c+\frac{2l}{5}\omega\\
&=\frac{96l^2}{D}v_0
+\frac{2l}{5}\frac{400x}{D}v_0\\
&=\boxed{\frac{96l^2+160lx}{D}v_0}.
\end{aligned}\]

速度代数值的正负表示相对于入射方向的正、反向。上述顺序避免了直接求解三元二次方程，也使每一次代回都有明确来源。

\emph{【物理校验】} 当 \(x=0\) 时 \(\omega=0\)，碰撞退化为质量 \(80m\)
与 \(5m\) 的一维完全弹性碰撞：

\[v=\frac{80m-5m}{80m+5m}v_0=\frac{15}{17}v_0,\]

而上式在 \(x=0\) 时给出 \(45/51=15/17\)，完全一致。

\textbf{（2）}

碰后入射小球与哑铃质心的相对水平速度为
\(v-V_c\)。要发生第二次碰撞，小球既不能始终落在哑铃后方，也不能在杆尚未转到可相交方位前远离。

设碰后经过时间 \(t\)，杆相对初始竖直方向转过

\[\theta=\omega t.\]

以哑铃质心为原点，杆上距质心有向距离为 \(\xi\) 的点，其位置为

\[X=V_ct+\xi\sin\theta,\qquad
Y=\xi\cos\theta.\]

入射小球的轨迹为

\[X_b=vt,\qquad Y_b=x.\]

两者相交要求

\[\xi=\frac{x}{\cos\theta},\qquad
(v-V_c)t=x\tan\theta. \tag{4}\]

同时接触点必须落在真实杆段上：

\[-\frac{3l}{5}\leq\xi\leq\frac{2l}{5}. \tag{5}\]

令

\[y=\frac{x}{l},\qquad
k(y)=\frac{V_c-v}{x\omega}
=\frac{51-200y^2}{400y^2}.\]

下面说明式（6）的来源。把 \(t=\theta/\omega\) 代入式（4）的第二式：

\[(v-V_c)\frac{\theta}{\omega}=x\tan\theta.\]

两边除以 \(x\)，再利用 \(v-V_c=-(V_c-v)\)：

\[-\frac{V_c-v}{x\omega}\theta=\tan\theta,\]

因此

\[k(y)\theta+\tan\theta=0.\]

在 \(0<\theta<\pi/2\)
内，\(k(y)\theta>0\)、\(\tan\theta>0\)，两项不可能相加为零；所以第一次可能的再相交必须发生在第二象限
\(\theta\in(\pi/2,\pi)\)。式（4）于是化为

\[F_y(\theta)\equiv k(y)\theta+\tan\theta=0. \tag{6}\]

还要确认这个区间内确实只有一个候选交点。第一次碰撞点满足
\(0<y<2/5\)，因而

\[k(y)=\frac{51-200y^2}{400y^2}>0.\]

当 \(\theta\to(\pi/2)^+\) 时，\(\tan\theta\to-\infty\)，所以
\(F_y(\theta)\to-\infty\)；当 \(\theta\to\pi^-\)
时，\(\tan\theta\to0^-\)，所以

\[F_y(\theta)\to k(y)\pi>0.\]

又有

\[\frac{\partial F_y}{\partial\theta}=k(y)+\sec^2\theta>0,\]

故 \(F_y\)
在第二象限严格单调递增，并且恰有一个零点。这个零点对应``无限直线''首次再次相交的转角，后面只需判断交点是否落在有限杆段上。

在该区间内 \(\cos\theta<0\)，因此
\(\xi=x/\cos\theta<0\)，交点位于杆的质心下方，即朝 \(2m\)
端。真实杆段的下界要求

\[\xi=\frac{x}{\cos\theta}\geq-\frac{3l}{5}.\]

因为 \(\cos\theta<0\)，两边乘以 \(\cos\theta\) 时不等号反向：

\[x\leq-\frac{3l}{5}\cos\theta.\]

再除以 \(3l/5>0\)，得到交点没有越过 \(2m\) 端点的条件

\[-\cos\theta\geq\frac{5y}{3}. \tag{7}\]

临界状态同时满足式（6）、（7）的等号。令

\[\theta_e=\arccos\left(-\frac{5y}{3}\right),\]

则临界参数 \(y_c\) 满足一元方程

\[\boxed{
\frac{51-200y_c^2}{400y_c^2}
\arccos\left(-\frac{5y_c}{3}\right)
+\tan\!\left[\arccos\left(-\frac{5y_c}{3}\right)\right]=0
}. \tag{8}\]

数值求根得到

\[y_c=0.3204287981\ldots.\]

下面说明为什么临界根的右侧而不是左侧对应真实碰撞。令

\[\theta_e(y)=\arccos\left(-\frac{5y}{3}\right),\]

它表示交点恰好位于 \(2m\) 端时的转角。由于 \(-\cos\theta\) 在第二象限随
\(\theta\) 单调增大，交点留在杆段内等价于

\[\theta_*(y)\geq\theta_e(y),\]

其中 \(\theta_*(y)\) 是 \(F_y(\theta)=0\) 的唯一根。又因为 \(F_y\) 随
\(\theta\) 严格递增，上式进一步等价于

\[F_y\!\left(\theta_e(y)\right)\leq0.\]

把 \(\theta_e(y)\) 代入所得的等号正是式（8）。对 \(0<y<2/5\)
作一元数值检查可见，该函数只在 \(y_c\) 处过零，并在 \(y>y_c\)
时为负，因此允许端点接触时，第二次碰撞要求 \(y\ge y_c\)。

另一方面，第一次碰撞点位于质心上方的杆段内要求

\[0<y\le\frac25.\]

因此发生第二次有限冲量碰撞的条件为

\[\boxed{0.3204287981\ldots\le\frac{x}{l}\le\frac25}.\]

数值 \(0.32\) 只是保留两位小数的近似值，精确下限并不等于
\(8/25\)。下限对应第二次交点恰在 \(2m\) 端点，上限对应第一次碰撞点到达
\(3m\)
端点。理想质点哑铃允许碰端点，因此包含两端；若额外规定只能碰杆的内部，才改为开区间。端点接触本身并不意味着冲量为零。

\emph{【数学技巧：先判根存在，再求临界参数】} 第（2）问无需求出每个
\(y\)
对应的碰撞时刻。只需令交点恰好落在端点，联立相交方程与端点条件，即可得到参数边界。

\textbf{（3）}

设两次碰撞间隔为 \(t\)，并记杆的转角为

\[\theta=\omega t.\]

由第（2）问的相交条件

\[(v-V_c)t=x\tan\theta,\]

在允许第二次相交的范围内有 \(V_c>v\)，且首次再相交发生在第二象限
\(\pi/2<\theta<\pi\)。因此

\[\tan\theta=-\frac{(V_c-v)t}{x}.\]

令

\[\alpha=\arctan\frac{(V_c-v)t}{x},\qquad 0<\alpha<\frac\pi2.\]

第二象限中正切为负，所以正确的角度分支不是 \(-\alpha\)，而是

\[\theta=\pi-\alpha.\]

也就是

\[\boxed{\omega t=\pi-\arctan\frac{(V_c-v)t}{x}}. \tag{9}\]

若只要求一元方程，直接在正切式中代入

\[V_c-v=\frac{51l^2-200x^2}{51l^2+200x^2}v_0,\]

以及

\[\omega=\frac{400xv_0}{51l^2+200x^2},\]

即可写成仅含 \(t\) 的方程

\[\boxed{
\frac{v_0t}{x}\frac{51l^2-200x^2}{51l^2+200x^2}
=-\tan\left(\frac{400xv_0t}{51l^2+200x^2}\right)}.
\]

求出候选根后还要依次做三项筛选：\(t>0\)；转角位于首次允许区间；由
\(\xi=x/\cos(\omega t)\) 算出的接触位置属于
\([-3l/5,2l/5]\)。满足这些条件的最小正根才是两次碰撞之间的时间。

\textbf{端点不是相切的检验。} 取杆的单位方向
\(\boldsymbol e=(\sin\theta,\cos\theta)\)，顺时针法向
\(\boldsymbol n=(\cos\theta,-\sin\theta)\)。交点的杆上坐标为
\(\xi=x/\cos\theta\)，小球对该杆点的相对法向速度为

\[v_n=(v-V_c)\cos\theta-\xi\omega.\]

下限交点 \(\xi=-3l/5\)。代入临界数值
\(x/l=0.3204287981\)、\(\theta=\arccos[-5x/(3l)]\)，得到
\(v_n\approx1.30248v_0\)，非零。该接触具有有限相对速度，不能因``恰在端点''而删除。

\newpage

\section{第 3 题}\label{ux7b2c-3-ux9898}

一个均匀电阻圆环被若干等分点分割，圆环每一小段的电阻均为
\(2r\)，各等分点通过辐条与圆心 \(O\) 相连。

\begin{enumerate}
\def\labelenumi{\arabic{enumi}.}
\tightlist
\item
  如图（a），共有三根辐条，每根辐条的电阻均为 \(4r\)，求 \(A\)、\(O\)
  间的等效电阻 \(R_{AO1}\)。
\item
  如图（b），共有 \(2N\) 根辐条，\(N>2\)。位于同一直径上的
  \(OA\)、\(OB\) 两根辐条电阻为 \(r\)，其余辐条电阻均为 \(4r\)，求
  \(R_{AO2}\)。
\item
  如图（c），共有 \(N\) 根辐条，\(N>2\)，每根辐条的电阻均为 \(4r\)，求
  \(R_{AO3}\)。
\end{enumerate}

\begin{figure}[H]
\centering
\includegraphics[width=0.9\linewidth,height=\textheight,keepaspectratio,alt={图 3-1　三种具有旋转或镜面对称性的辐条---圆环电阻网络。}]{../figures/question-03.pdf}
\caption*{图 3-1　三种具有旋转或镜面对称性的辐条---圆环电阻网络。}
\end{figure}

\subsection{一、题型定位}\label{ux4e00ux9898ux578bux5b9aux4f4d-2}

先利用对称性合并等势节点，再写电阻递推。环形网络的末问需列节点方程；差分方程的边界条件决定最终等效电阻。

\subsection{二、知识点}\label{ux4e8cux77e5ux8bc6ux70b9-2}

\begin{itemize}
\tightlist
\item
  对称电路中的等势点判断。
\item
  星形网络、环形网络与梯形网络的等效化。
\item
  节点电压法与二阶常系数差分方程。
\end{itemize}

\subsection{三、物理图像}\label{ux4e09ux7269ux7406ux56feux50cf-2}

给 \(A\)、\(O\)
两端加测试电压后，电势从端口沿圆环向两侧传播。对称节点等势，可折叠网络；一般情形中，相邻节点电势按固定递推关系衰减。

\subsection{四、解题入口}\label{ux56dbux89e3ux9898ux5165ux53e3-2}

在端口 \(AO\)
接测试电压，先找等势点。能靠对称性折叠就先折叠；不能彻底折叠时，对任一普通节点列基尔霍夫电流方程。

本题的难点不在并联、串联公式，而在``如何把复杂网络变简单''。前三步依次是：先给端口加测试电压；再用对称性寻找等势点；最后才合并电阻。第（3）问无法完全靠直观化简时，再升级为节点电压法，把电路变成一个数列递推问题：

\[\text{几何对称}\longrightarrow\text{等势合并}\longrightarrow\text{节点递推}.\]

判断节点是否等势时，还要检查外接电源是否保持同一对称性；图形的几何对称本身并不保证节点等势。

\subsection{五、解析}\label{ux4e94ux89e3ux6790-2}

图 3-2 把关于端口轴 \(AO\)
镜面对称的环形电阻网络折叠成梯形网络，并标出节点电压递推
\(2U_{i-1}-5U_i+2U_{i+1}=0\)。第（1）问先以等势节点完成有限网络化简；第（2）、（3）问使用折叠后的梯形链和递推式处理多级或
\(N\) 辐条情形。

\begin{figure}[H]
\centering
\includegraphics[width=0.569\linewidth,height=\textheight,keepaspectratio,alt={图 3-2　镜面对称折叠得到的梯形网络与节点电压递推。}]{../figures/solution-03.pdf}
\caption*{图 3-2　镜面对称折叠得到的梯形网络与节点电压递推。}
\end{figure}

\textbf{（1）}

图 3-2 所用折叠成立的前提是测试电源也关于 \(AO\)
对称。只有在这个边界条件下，镜像位置的节点才等势，相连电阻才可判为无电流。

为了说明每一步等效的依据，给 \(A\)、\(O\) 两端接入一个测试电源，设

\[V_A=U,\qquad V_O=0.\]

把圆周上另外两个节点分别记为 \(B\)、\(C\)。整个电路关于直线 \(AO\)
镜面对称，而测试电源的接法也关于 \(AO\) 对称，因此交换 \(B\)、\(C\)
后电路与边界条件均不改变。唯一的节点电势解必须满足

\[V_B=V_C.\]

所以连接 \(B\)、\(C\) 的那一段圆环电阻两端电压为零，其电流为

\[I_{BC}=\frac{V_B-V_C}{2r}=0.\]

这根电阻可以从等效电路中删去，同时可把 \(B\)、\(C\) 合并成一个等势节点
\(D\)。合并后，从 \(A\) 到 \(D\) 有两根阻值均为 \(2r\)
的圆环电阻并联，因此

\[2r\parallel2r=r,\]

从 \(D\) 到 \(O\) 有原来 \(BO\)、\(CO\) 两根 \(4r\) 辐条并联，因此

\[4r\parallel4r=2r.\]

于是经过节点 \(D\) 的通路为串联：

\[R_{ADO}=r+2r=3r.\]

除此之外，\(A\) 与 \(O\) 之间还有一根直接连接的 \(4r\) 辐条。因此端口
\(AO\) 实际看到的是 \(4r\) 与 \(3r\) 并联：

\[\begin{aligned}
\frac1{R_{AO1}}
&=\frac1{4r}+\frac1{3r}\\
&=\frac{3+4}{12r}\\
&=\frac7{12r}.
\end{aligned}\]

\[\boxed{R_{AO1}=4r\parallel3r=\frac{12}{7}r}.\]

\textbf{（2）}

仍然在 \(A\)、\(O\) 之间加测试电压。直径 \(AB\)
上方和下方的电路互为镜像，因此上、下圆周对应节点等势，可以逐对合并。

合并时要注意两类电阻的变化：

\begin{enumerate}
\def\labelenumi{\arabic{enumi}.}
\tightlist
\item
  上、下圆周对应的两段 \(2r\) 电阻连接同一对等势节点，故二者并联为
  \[2r\parallel2r=r.\]
\item
  除 \(A\)、\(B\) 以外，每一对普通辐条均为两根 \(4r\) 并联，故合并后为
  \[4r\parallel4r=2r.\]
\end{enumerate}

折叠后的电路是一串有限梯形单元。为避免一下子看完整个网络，定义 \(X_j\)
为``从第 \(j\) 个折叠节点向右看，直到 \(O\) 的等效电阻''。从最右端 \(B\)
开始：\(BO\) 是题目中特别给出的电阻 \(r\)，所以

\[X_N=r.\]

向左退一级时，该节点有两条到 \(O\) 的通路：一条是本节点的竖直电阻
\(2r\)；另一条是先经过横向电阻 \(r\)，再经过右侧等效电阻
\(X_{j+1}\)。因此递推关系是

\[X_j=2r\parallel(r+X_{j+1}). \tag{1}\]

若某一级 \(X_{j+1}=r\)，则

\[2r\parallel(r+r)=r.\]

由于末端 \(X_N=r\)，根据上述递推可反向归纳得到

\[X_{N-1}=X_{N-2}=\cdots=X_1=r.\]

从 \(A\) 沿圆环进入梯形网络时，还要先经过最左端的一段横向电阻
\(r\)，所以除 \(OA\) 直连辐条以外的整条通路等效为

\[r+X_1=r+r=2r.\]

而 \(OA\) 本身是题目特别给出的电阻 \(r\)。最终端口电阻为

\[\begin{aligned}
R_{AO2}
&=r\parallel2r\\
&=\frac{r\cdot2r}{r+2r}\\
&=\frac23r.
\end{aligned}\]

\[\boxed{R_{AO2}=r\parallel2r=\frac23r}.\]

结果与 \(N\) 无关，反映了特殊端点阻值和镜面对称共同造成的电势锁定。

\textbf{（3）}

第三问所有辐条均为 \(4r\)。设圆心 \(O\) 为零电势，端口 \(A\) 的电势为
\(U\)。沿圆环依次给节点编号

\[0,1,2,\ldots,N-1,\]

其中节点 \(0\) 就是 \(A\)。为处理周期闭合，约定节点 \(N\) 与节点 \(0\)
是同一个物理节点。令第 \(i\) 个圆周节点的电势为 \(U_i\)，于是

\[U_0=U_N=U.\]

现在选取一个非端口节点 \(i\)。从该节点流出的电流有三股：

\begin{itemize}
\tightlist
\item
  经左侧圆环电阻 \(2r\) 流向节点 \(i-1\)；
\item
  经右侧圆环电阻 \(2r\) 流向节点 \(i+1\)；
\item
  经辐条电阻 \(4r\) 流向零电势的圆心 \(O\)。
\end{itemize}

根据基尔霍夫电流定律，流出电流代数和为零：

\[\frac{U_i-U_{i-1}}{2r}
+\frac{U_i-U_{i+1}}{2r}
+\frac{U_i-0}{4r}=0.\]

方程两边乘以 \(4r\)：

\[2(U_i-U_{i-1})+2(U_i-U_{i+1})+U_i=0.\]

展开并整理：

\[5U_i-2U_{i-1}-2U_{i+1}=0,\]

等价地写成

\[2U_{i-1}-5U_i+2U_{i+1}=0. \tag{1}\]

这是一个二阶常系数差分方程。仿照常系数微分方程的指数试探法，令

\[U_i=\lambda^i.\]

代入式（1）：

\[2\lambda^{i-1}-5\lambda^i+2\lambda^{i+1}=0.\]

当 \(\lambda\neq0\) 时，两边除以 \(\lambda^{i-1}\)：

\[2-5\lambda+2\lambda^2=0.\]

把它写成通常的降幂顺序：

\[2\lambda^2-5\lambda+2=0
=(2\lambda-1)(\lambda-2).\]

其根为 \(\lambda=2\) 与 \(\lambda=\frac12\)，故

\[U_i=A2^i+B2^{-i}. \tag{2}\]

令端口电压为 \(U_0=U\)，并按网络的周期闭合条件连接第 \(N\)
个节点与端点。通解中的两项分别描述沿圆环顺、逆方向的指数衰减。下面把两个边界条件逐步代入。由
\(U_0=U\)，

\[A+B=U. \tag{3}\]

由 \(U_N=U\)，

\[A2^N+B2^{-N}=U. \tag{4}\]

令 \(q=2^N\)，式（4）成为

\[qA+\frac Bq=U.\]

两边乘以 \(q\)：

\[q^2A+B=qU. \tag{5}\]

由式（3）得 \(B=U-A\)，代入式（5）：

\[q^2A+U-A=qU.\]

因此

\[A(q^2-1)=U(q-1).\]

利用 \(q^2-1=(q-1)(q+1)\)，得到

\[A=\frac{U}{2^N+1},\qquad
B=\frac{2^NU}{2^N+1}.\]

接下来求与端点 \(A\) 相邻的节点电势 \(U_1\)：

\[\begin{aligned}
U_1
&=A2^1+B2^{-1}\\
&=2A+\frac B2\\
&=\frac{2U}{2^N+1}
+\frac{2^NU}{2(2^N+1)}\\
&=\frac{(4+2^N)U}{2(2^N+1)}.
\end{aligned}\]

端点 \(A\) 的另一侧相邻节点是 \(N-1\)。将 \(i=N-1\) 代入通解并使用上面的
\(A\)、\(B\)：

\[\begin{aligned}
U_{N-1}
&=A2^{N-1}+B2^{-(N-1)}\\
&=\frac{2^{N-1}U}{2^N+1}
+\frac{2^NU}{2^N+1}\,2^{-(N-1)}\\
&=\frac{(2^N+4)U}{2(2^N+1)}\\
&=U_1.
\end{aligned}\]

这一步把``左右两侧相邻节点电势相同''从几何直觉落实为通解的直接结果。

端点流入网络的总电流由直达中心的辐条电流与左右两段圆环电流组成：

\[I=\frac{U}{4r}
+2\frac{U-U_1}{2r}. \tag{6}\]

式（6）中第一项是 \(A\) 点经本身辐条流向圆心的电流。后两项是 \(A\)
点分别经顺时针、逆时针两段圆环流出的电流；由于两侧相邻节点电势相同，所以出现因子
\(2\)。

把 \(U_1\) 代入：

\[\begin{aligned}
I
&=\frac{U}{4r}+\frac{U-U_1}{r}\\
&=\frac{U}{4r}
+\frac Ur\left[1-\frac{2^N+4}{2(2^N+1)}\right]\\
&=\frac{U}{4r}
+\frac Ur\frac{2^{N+1}+2-2^N-4}{2(2^N+1)}\\
&=\frac{U}{4r}
+\frac Ur\frac{2^N-2}{2(2^N+1)}\\
&=\frac{3(2^N-1)}{4(2^N+1)}\frac Ur.
\end{aligned}\]

因此

\[\boxed{
R_{AO3}=\frac{4(2^N+1)}{3(2^N-1)}\,r
}.\]

当 \(N\to\infty\) 时，

\[R_{AO3}\longrightarrow \frac43r,\]

可作为结果的量级检验。

还可检查最小允许值 \(N=3\)：此时
\(R_{AO3}=4(8+1)r/[3(8-1)]=12r/7\)，与第（1）问的三辐条结果一致，边界条件检查通过。

图 3-2 上图的下导线都是节点 \(O\)，上端右侧是节点
\(B\)。每个锯齿代表一个电阻，不能另画导线跨接其两端。下图把第（3）问的圆环展开为周期链：\(A_0=A_N=A\)，横向电阻为
\(2r\)，每一辐条为
\(4r\)。对普通节点的三股电流求和，直接得到式（1）；边界必须用
\(U_N=U_0\)，不能把右端设为接地。

\newpage

\section{第 4 题}\label{ux7b2c-4-ux9898}

一架无机械转动结构的主动动力飞机总质量为
\(m\)，其动力来自高压组件产生的离子流。发射电极与收集电极相距
\(l\)，两电极间电压为
\(U\)，其间电场可视为匀强电场。空气中的氮气被电离为
\(\mathrm{N_2^+}\)，离子从发射极由静止出发。元电荷为
\(e\)，阿伏伽德罗常数为 \(N_A\)，氮气摩尔质量为 \(M_{\mathrm{N_2}}\)。

\begin{enumerate}
\def\labelenumi{\arabic{enumi}.}
\tightlist
\item
  忽略离子与空气分子的碰撞，求离子到达收集电极前的速度 \(v_1\)。
\item
  设离子的平均自由程为 \(\lambda\)，每次碰撞间位移均为
  \(\lambda\)，且碰撞后离子速度归零。求每个离子从发射极到收集极的过程中传递给空气的总动量
  \(\delta p\)。
\item
  飞机所受阻力为 \(f=-\alpha v\)。若共有 \(N\) 组电极，每组电流均为
  \(I\)，求飞机稳定运动时的速度 \(v_2\)。
\item
  定义推功比为推力与输入功率之比，求 \(\eta\)。
\item
  飞机平稳飞行需满足 \[nmg\leq\frac12\rho v^2SC_L,\] 其中 \(n\)
  为负荷系数，\(\rho\) 为空气密度，\(S\) 为机翼面积，\(C_L\)
  为升力系数。求稳定飞行所需的最小单组电流 \(I_0\)。
\end{enumerate}

\begin{figure}[H]
\centering
\includegraphics[width=0.9\linewidth,height=\textheight,keepaspectratio,alt={图 4-1　离子在发射电极与收集电极之间受电场加速并形成离子风。}]{../figures/question-04.pdf}
\caption*{图 4-1　离子在发射电极与收集电极之间受电场加速并形成离子风。}
\end{figure}

\subsection{一、题型定位}\label{ux4e00ux9898ux578bux5b9aux4f4d-3}

从单个离子在自由程内获得的动量算起，用电流换算单位时间碰撞次数，求空气推力。再按阻力模型求速度、功率和起飞所需电流。

\subsection{二、知识点}\label{ux4e8cux77e5ux8bc6ux70b9-3}

\begin{itemize}
\tightlist
\item
  电场力做功与带电粒子的加速。
\item
  平均自由程模型中的分段加速和动量累积。
\item
  电流与单位时间离子数的换算。
\item
  推力---阻力平衡及升力约束。
\end{itemize}

\subsection{三、物理图像}\label{ux4e09ux7269ux7406ux56feux50cf-3}

每个离子在一个平均自由程内从静止加速，碰撞时把动量交给中性空气，随后重新从静止加速。许多离子连续重复这一过程，形成稳定的动量流和宏观推力。

\subsection{四、解题入口}\label{ux56dbux89e3ux9898ux5165ux53e3-3}

始终从一个离子开始：先把摩尔质量换成单粒子质量，再用一个自由程内的电场功求碰前动量。之后才乘碰撞次数和单位时间离子数。

先沿着``一个离子''的经历走完整条链条，再扩大到整架飞机。一个离子在自由程内获得动能，碰撞时把动量交给空气；每秒通过的离子数由电流给出；单位时间传递的动量就是推力；推力再与阻力、升力条件连接。全题只有一条主线：

\[\text{电场做功}\to\text{单次碰撞动量}\to\text{单位时间动量}\to\text{飞机受力}.\]

高一学生最容易把``动量''\,``动量变化率''和``力''混为一谈，讲解时应在每一步都写出单位。

\subsection{五、解析}\label{ux4e94ux89e3ux6790-3}

图 4-2 把离子在每个平均自由程 \(\lambda\)
内的运动画成重复的锯齿速度图：离子由近似静止受电场加速，碰撞时速度近似归零，该段获得的动量全部传给空气。第（1）问求无碰撞末速；第（2）问由单段加速求每次碰撞的动量；第（3）问用单位时间的离子数乘单离子总动量转移率。

\begin{figure}[H]
\centering
\includegraphics[width=0.9\linewidth,height=\textheight,keepaspectratio,alt={图 4-2　离子在相邻碰撞间加速并逐段向空气传递动量。}]{../figures/solution-04.pdf}
\caption*{图 4-2　离子在相邻碰撞间加速并逐段向空气传递动量。}
\end{figure}

\textbf{（1）}

图 4-2
的每一条斜线对应匀强电场中的匀加速，竖直回落对应碰撞后有序漂移速度近似归零。第一问相当于只有一段、长度取完整极间距
\(l\) 的极限。

氮气的摩尔质量 \(M_{\mathrm{N_2}}\) 是一摩尔氮气分子的质量。一摩尔含有
\(N_A\) 个分子，所以单个氮气分子、也就是单个 \(\mathrm{N_2^+}\)
离子的质量为

\[m_i=\frac{M_{\mathrm{N_2}}}{N_A}.\]

离子带一个正元电荷
\(+e\)。它从发射电极运动到收集电极，跨越的电势差大小为
\(U\)，电场力所做的功为

\[W_e=eU.\]

离子从静止开始，且本问忽略碰撞，所以电场功全部转化为离子动能：

\[eU=\frac12m_iv_1^2.\]

把 \(m_i=M_{\mathrm{N_2}}/N_A\) 代入：

\[eU=\frac12\frac{M_{\mathrm{N_2}}}{N_A}v_1^2.\]

两边乘以 \(2N_A/M_{\mathrm{N_2}}\)：

\[v_1^2=\frac{2eUN_A}{M_{\mathrm{N_2}}}.\]

速度取正的大小，因此

\[\boxed{v_1=\sqrt{\frac{2eUN_A}{M_{\mathrm{N_2}}}}}.\]

也可由匀强电场运动学验证。离子加速度

\[a=\frac{eE}{m_i}=\frac{eUN_A}{lM_{\mathrm{N_2}}},\]

从静止运动距离 \(l\) 后

\[v_1^2=2al=\frac{2eUN_A}{M_{\mathrm{N_2}}},\]

与能量法结果一致。能量法不需要显式求飞行时间，因而更简洁。

\textbf{（2）}

两电极间电场均匀，场强大小为

\[E=\frac Ul.\]

在一个平均自由程内，离子从静止出发并沿电场方向运动距离
\(\lambda\)。这一小段内电场力做功为

\[eE\lambda=\frac{eU\lambda}{l}.\]

设离子在下一次碰撞前的速度为 \(v_\lambda\)。由动能定理，

\[\frac12m_iv_\lambda^2=\frac{eU\lambda}{l}.\]

故

\[v_\lambda=\sqrt{\frac{2eU\lambda}{m_il}}.\]

碰撞前离子动量大小为

\[\begin{aligned}
p_\lambda
&=m_iv_\lambda\\
&=m_i\sqrt{\frac{2eU\lambda}{m_il}}\\
&=\sqrt{\frac{2m_ieU\lambda}{l}}.
\end{aligned}\]

题设假定每次碰撞后离子速度归零，因此每个自由程都重复相同的``从静止加速---碰撞归零''过程。电极间总距离为
\(l\)，每段长度为 \(\lambda\)，所以碰撞段数近似为

\[n_c=\frac l\lambda.\]

每次碰撞时，离子损失动量
\(p_\lambda\)；根据动量守恒，空气获得同样大小的动量。故单个离子在全过程传递给空气的总动量为

\[\begin{aligned}
\delta p
&=n_cp_\lambda\\
&=\frac l\lambda
\sqrt{\frac{2m_ieU\lambda}{l}}\\
&=\sqrt{\frac{l^2}{\lambda^2}
\frac{2m_ieU\lambda}{l}}\\
&=\sqrt{\frac{2m_ieUl}{\lambda}}.
\end{aligned}\]

最后代入 \(m_i=M_{\mathrm{N_2}}/N_A\)：

\[\delta p=\frac{l}{\lambda}p_\lambda
=\boxed{\sqrt{\frac{2eUM_{\mathrm{N_2}}l}{\lambda N_A}}}.\]

这里默认 \(l/\lambda\) 足够大，可忽略最后一个不足 \(\lambda\)
的余段。若必须保留离散性，应令完整碰撞次数为
\(\lfloor l/\lambda\rfloor\)，并另行处理末段；竞赛模型采用连续近似即可。

注意总动量不是无碰撞末动量的 \(l/\lambda\) 倍。每段碰撞前的动量只按
\(\sqrt{\lambda}\) 增长，而碰撞次数按 \(1/\lambda\) 增长，故

\[\delta p\propto\lambda^{-1/2}.\]

\textbf{（3）}

电流的定义是单位时间通过截面的电荷量：

\[I=\frac{\Delta Q}{\Delta t}.\]

每个离子带电量 \(e\)。若每秒通过的离子数为 \(\dot N_i\)，则

\[I=e\dot N_i,\]

所以一组电极单位时间通过的离子数为

\[\dot N_i=\frac Ie.\]

每个离子给空气传递动量
\(\delta p\)，故一组电极单位时间传递的总动量，即推力大小为

\[F_1=\dot N_i\delta p=\frac Ie\delta p.\]

共有 \(N\) 组相同电极，所以总推力为

\[F=N\frac{I}{e}\delta p.\]

飞机速度为 \(v\) 时，阻力大小为
\(\alpha v\)，方向与速度相反。稳定运动意味着加速度为零，即水平方向合力为零：

\[F-\alpha v_2=0.\]

因此

\[v_2=\frac F\alpha
=\frac{NI}{\alpha e}\delta p.\]

代入第（2）问的 \(\delta p\)：

\[\boxed{
v_2=\frac{NI}{\alpha e}
\sqrt{\frac{2eUM_{\mathrm{N_2}}l}{\lambda N_A}}
}.\]

推力方向与离子加速方向相反，这是空气获得动量后对机体的反作用。上式只写速度大小；若以离子流方向为正，飞机速度应带相反符号。

\textbf{（4）}

每组电极两端电压为 \(U\)、电流为 \(I\)，其输入电功率为 \(UI\)。\(N\)
组电极的总输入功率为

\[P_{\mathrm{in}}=NUI.\]

总推力为

\[F=\frac{NI}{e}\delta p.\]

根据推功比定义，

\[\begin{aligned}
\eta
&=\frac{F}{P_{\mathrm{in}}}\\
&=\frac{(NI/e)\delta p}{NIU}\\
&=\frac{\delta p}{eU}.
\end{aligned}\]

这里 \(N\) 与 \(I\)
同时约去，说明在这个理想模型中，增加电极组数会同比增加推力和输入功率，却不改变推功比。代入
\(\delta p\)：

\[\eta=\frac{F}{P}
=\frac{\delta p}{eU}
=\boxed{\sqrt{\frac{2M_{\mathrm{N_2}}l}{eU\lambda N_A}}}.\]

量纲上，\(\eta=F/P\) 的单位为

\[\frac{\mathrm N}{\mathrm W}
=\frac{\mathrm s}{\mathrm m},\]

即速度的倒数，并非无量纲效率。名称``推功比''强调单位输入功率可产生的推力。

\textbf{（5）}

飞机刚好能够满足稳定飞行要求时，升力条件取等号：

\[nmg=\frac12\rho v_{\min}^2SC_L.\]

两边乘以 \(2/(\rho SC_L)\)：

\[v_{\min}^2=\frac{2nmg}{\rho SC_L}.\]

取正的速度大小：

\[v_{\min}=\sqrt{\frac{2nmg}{\rho SC_L}}.\]

飞机的稳定速度 \(v_2\) 与每组电流 \(I\)
成正比。要使电流最小，就令稳定速度恰好等于最低允许速度：

\[v_2=v_{\min}.\]

代入第（3）问的表达式：

\[\frac{NI_0}{\alpha e}
\sqrt{\frac{2eUM_{\mathrm{N_2}}l}{\lambda N_A}}
=\sqrt{\frac{2nmg}{\rho SC_L}}.\]

先把根号因子移到右边：

\[I_0=\frac{\alpha e}{N}
\frac{\sqrt{2nmg/(\rho SC_L)}}
{\sqrt{2eUM_{\mathrm{N_2}}l/(\lambda N_A)}}.\]

利用 \(\sqrt a/\sqrt b=\sqrt{a/b}\)：

\[\begin{aligned}
I_0
&=\frac{\alpha e}{N}
\sqrt{
\frac{2nmg}{\rho SC_L}
\frac{\lambda N_A}{2eUM_{\mathrm{N_2}}l}
}\\
&=\frac{\alpha e}{N}
\sqrt{\frac{nmg\lambda N_A}
{\rho SC_L\,eUM_{\mathrm{N_2}}l}}.
\end{aligned}\]

因此

\[\boxed{
I_0=\frac{\alpha e}{N}
\sqrt{\frac{nmg\lambda N_A}
{\rho SC_L\,eU M_{\mathrm{N_2}}l}}
}.\]

若题目把 \(I_0\) 定义为所有电极的总电流，则应将结果中的 \(1/N\)
去掉；下文按``每组电流''理解。

图 4-2 的横轴是路程，故每段速度满足
\(v^2=2(eU/ml)(s-s_i)\)，曲线为平方根形；若改画速度---时间图才是直线。推力估算把
\(l/\lambda\) 当作很大的碰撞次数，忽略最后不足一段平均自由程的动量。若
\(l/\lambda\) 不是整数，真正传给空气的完整碰撞段数应取
\(\lfloor l/\lambda\rfloor\)；末段离子到达收集极的剩余动量主要传给电极，不能再算作离子风的空气动量。

\newpage

\section{第 5 题}\label{ux7b2c-5-ux9898}

交流电路如图。已知 \(C_0\)、\(L_0\)、\(R\)、\(C_1\)，交流电源电压为

\[U(t)=u\cos\omega t.\]

\begin{enumerate}
\def\labelenumi{\arabic{enumi}.}
\tightlist
\item
  记前级 \(C_0\)、\(L_0\)、\(R\) 串联部分两端的电压为
  \(U_0\)。若要求总电流 \(I_0\) 与 \(U_0\) 同相，求圆频率 \(\omega\)
  应满足的条件。
\item
  在第（1）问条件下，若还要求支路电流 \(I_1\)、\(I_2\) 相对于 \(I_0\)
  的相位对称，即二者分别超前、滞后 \(I_0\) 相同的角度，求 \(L_1\)。
\item
  求 \(I_0(t)\)、\(I_1(t)\)、\(I_2(t)\)，并求电路总平均功率。
\end{enumerate}

\begin{figure}[H]
\centering
\includegraphics[width=0.9\linewidth,height=\textheight,keepaspectratio,alt={图 5-1　电路结构：L\_0、C\_0、R 串联后接入 RL\_1 与 RC\_1 两并联支路，U\_0 为前级串联段电压。}]{../figures/question-05.pdf}
\caption*{图 5-1　电路结构：\(L_0\)、\(C_0\)、\(R\) 串联后接入 \(RL_1\)
与 \(RC_1\) 两并联支路，\(U_0\) 为前级串联段电压。}
\end{figure}

\subsection{一、题型定位}\label{ux4e00ux9898ux578bux5b9aux4f4d-4}

用相量计算正弦交流稳态。串联谐振和共轭阻抗简化电流叠加；求功率时仍须分清峰值与有效值。

\subsection{二、知识点}\label{ux4e8cux77e5ux8bc6ux70b9-4}

\begin{itemize}
\tightlist
\item
  正弦交流电的复阻抗表示。
\item
  串联谐振与并联支路相位对称。
\item
  相量叠加及有效值平均功率。
\end{itemize}

\subsection{三、物理图像}\label{ux4e09ux7269ux7406ux56feux50cf-4}

把各电流相量看成复平面上的箭头。电感支路电流位于实轴下方，电容支路电流位于实轴上方；当两箭头关于总电流对称时，它们的虚部抵消、实部相加。

\subsection{四、解题入口}\label{ux56dbux89e3ux9898ux5165ux53e3-4}

先准确识别 \(U_0\) 所跨的前级串联段，写出该段阻抗
\(Z_0\)。第（2）问再从两支路电流相位对称推出两支路阻抗互为共轭。

若学生尚未系统学习复数，不要先堆阻抗公式。先把每个相量看成一根以角速度
\(\omega\)
转动的箭头：电阻上的电压、电流同向；电感电流滞后；电容电流超前。随后再用复数把``箭头的长度与方向''压缩成一个量。解题顺序是

\[U_0\text{ 与 }I_0\text{ 同相}\to L_0C_0\text{ 串联谐振}\to
Z_1,Z_2\text{ 共轭}\to\text{求电流和功率}.\]

每算出一个相位，都要用``电感应滞后、电容应超前''检查一次方向。

\subsection{五、解析}\label{ux4e94ux89e3ux6790-4}

图 5-2 以并联端电压为实轴，画出 \(RL_1\) 支路电流 \(I_1\) 的滞后相角与
\(RC_1\) 支路电流 \(I_2\)
的超前相角；两相量关于实轴对称时，虚部抵消，总电流 \(I_0\)
沿实轴。第（2）问由共轭阻抗确定
\(L_1\)，第（3）问逐一投影相量得到三个时间函数和平均功率。

\begin{figure}[H]
\centering
\includegraphics[width=0.685\linewidth,height=\textheight,keepaspectratio,alt={图 5-2　两并联支路电流相量关于总电流对称。}]{../figures/solution-05.pdf}
\caption*{图 5-2　两并联支路电流相量关于总电流对称。}
\end{figure}

\textbf{（1）}

先明确本问所说的电压是图中前级串联段电压 \(U_0\)，而不是整个电源电压
\(U(t)\)。前级的复阻抗为

\[Z_0=R+j\omega L_0+\frac1{j\omega C_0}.\]

利用 \(1/j=-j\)，可写成

\[Z_0=R+j\left(\omega L_0-\frac1{\omega C_0}\right).\]

相量满足

\[\widetilde U_0=\widetilde I_0Z_0.\]

要使 \(I_0\) 与 \(U_0\) 同相，\(Z_0\) 必须是正实数，也就是其虚部等于零：

\[\omega L_0-\frac{1}{\omega C_0}=0.\]

逐步解这个方程：

\[\omega L_0=\frac1{\omega C_0},\]

\[\omega^2L_0C_0=1.\]

圆频率取正值，故

\[\boxed{\omega=\frac1{\sqrt{L_0C_0}}}.\]

在此频率下，电感与电容的电抗互相抵消，前级阻抗退化为
\(Z_0=R\)。并联支路的阻抗不包含在 \(U_0\)
的定义中，所以第（1）问此时不需要预先知道 \(L_1\)。

\textbf{（2）}

转到图
5-2。两支路电流关于总电流方向对称，等价于它们相对于并联端电压的相角大小相等、符号相反；因两支路电阻同为
\(R\)，这又等价于两个阻抗互为共轭。

上支路阻抗为

\[Z_1=R+j\omega L_1,\]

下支路阻抗为

\[Z_2=R-\frac{j}{\omega C_1}.\]

要使两支路电流关于总电流相位对称，两支路阻抗的虚部应大小相等、符号相反：

\[\omega L_1=\frac1{\omega C_1}.\]

因此

\[\boxed{L_1=\frac1{\omega^2C_1}=\frac{L_0C_0}{C_1}}.\]

更严格地，可设并联部分的端电压相量为 \(\widetilde U_p\)，并把
\(\widetilde U_p\)
选作相位参考，使它位于复平面的正实轴上。两支路电流相量为

\[\widetilde I_1=\frac{\widetilde U_p}{R+j\omega L_1},\qquad
\widetilde I_2=\frac{\widetilde U_p}{R-j/(\omega C_1)}.\]

要使 \(\widetilde I_1\)、\(\widetilde I_2\)
关于实轴对称，不仅要求相位绝对值相等，还要求两者模长相等。两支路电流的和
\(\widetilde I_0\)
随后也落在实轴上，所以这条实轴正是总电流的相位轴。由于两支路电阻均为
\(R\)，这两个要求都归结为

\[\omega L_1=\frac1{\omega C_1}.\]

因此该条件既是必要条件，也是充分条件。

也可以只比较模长来看到这一点。相位对称的两支路电流应有相同大小，而它们的端电压相同，因此必须满足

\[|Z_1|=|Z_2|.\]

平方后

\[R^2+(\omega L_1)^2
=R^2+\left(\frac1{\omega C_1}\right)^2.\]

两个电抗均取正的大小，所以仍得到
\(\omega L_1=1/(\omega C_1)\)。当两阻抗成为共轭复数时，两电流也是共轭相量，它们的和自然位于对称轴上，因而模长条件与相位条件同时满足。

\textbf{（3）}

定义

\[A=\frac{L_0C_0}{C_1^2},\qquad D=3R^2+A,\]

并令

\[\varphi=\arctan\frac{\sqrt A}{R}.\]

在已求条件下，令

\[X=\omega L_1=\frac1{\omega C_1}
=\frac{\sqrt{L_0C_0}}{C_1}=\sqrt A.\]

于是两支路阻抗为一对共轭复数：

\[Z_1=R+jX,\qquad Z_2=R-jX.\]

并联等效阻抗为

\[\begin{aligned}
Z_p&=\frac{Z_1Z_2}{Z_1+Z_2}\\
&=\frac{(R+jX)(R-jX)}{2R}\\
&=\frac{R^2+X^2}{2R}
=\frac{R^2+A}{2R}.
\end{aligned}\]

\emph{【数学技巧：共轭复数乘积】} 利用
\((R+jX)(R-jX)=R^2+X^2\)，并联部分的虚部自动消失。

由于 \(L_0\)、\(C_0\) 串联谐振，其总阻抗为零，电源看到的总阻抗为

\[Z_{\mathrm{tot}}
=R+Z_p
=R+\frac{R^2+A}{2R}
=\frac{3R^2+A}{2R}
=\frac D{2R}.\]

故总电流相量为

\[\widetilde I_0=\frac{u}{Z_{\mathrm{tot}}}
=\frac{2Ru}{D},\]

它是纯实数，说明 \(I_0\) 与电源电压同相。

并联部分的电压相量为

\[\widetilde U_p=\widetilde I_0Z_p
=\frac{u(R^2+A)}{D}.\]

所以

\[\begin{aligned}
\widetilde I_1
&=\frac{u(R^2+A)}{D(R+j\sqrt A)}\\
&=\frac{u(R^2+A)(R-j\sqrt A)}{D(R^2+A)}\\
&=\frac uD(R-j\sqrt A)\\
&=\frac{u\sqrt{R^2+A}}{D}e^{-j\varphi},\\
\widetilde I_2
&=\frac{u(R^2+A)}{D(R-j\sqrt A)}\\
&=\frac{u(R^2+A)(R+j\sqrt A)}{D(R^2+A)}\\
&=\frac uD(R+j\sqrt A)\\
&=\frac{u\sqrt{R^2+A}}{D}e^{+j\varphi}.
\end{aligned}\]

这里使用了

\[R\mp j\sqrt A=\sqrt{R^2+A}\,e^{\mp j\varphi},
\qquad \tan\varphi=\frac{\sqrt A}{R}.\]

把两个支路电流直接相加，还可核对节点电流定律：

\[\begin{aligned}
\widetilde I_1+\widetilde I_2
&=\frac uD(R-j\sqrt A)+\frac uD(R+j\sqrt A)\\
&=\frac{2Ru}{D}\\
&=\widetilde I_0.
\end{aligned}\]

虚部在相加时逐项抵消，因而总电流确实位于两支路电流的对称轴上。

因此电感支路 \(I_1\) 的相位为 \(-\varphi\)，即滞后；电容支路 \(I_2\)
的相位为 \(+\varphi\)，即超前。把相量 \(I_m e^{j\phi}\)
还原为瞬时量时，使用

\[\operatorname{Re}\!\left[I_me^{j\phi}e^{j\omega t}\right]
=I_m\cos(\omega t+\phi).\]

相量计算得到

\[\boxed{I_0(t)=\frac{2Ru}{D}\cos\omega t},\]

\[\boxed{I_1(t)=\frac{u\sqrt{R^2+A}}{D}\cos(\omega t-\varphi)},\]

\[\boxed{I_2(t)=\frac{u\sqrt{R^2+A}}{D}\cos(\omega t+\varphi)}.\]

电源电压峰值为 \(u\)，总电流峰值为 \(2Ru/D\)，且二者同相。因此平均功率为

\[\boxed{\overline P=\frac12u\frac{2Ru}{D}
=\frac{Ru^2}{3R^2+A}}.\]

还可用三个电阻的焦耳热独立校验。总平均功率为

\[\overline P
=\frac12R I_{0m}^2
+\frac12R I_{1m}^2
+\frac12R I_{2m}^2,\]

其中下标 \(m\) 表示峰值。代入

\[I_{0m}=\frac{2Ru}{D},\qquad
I_{1m}=I_{2m}=\frac{u\sqrt{R^2+A}}{D}\]

后，

\[\begin{aligned}
\overline P
&=\frac{R}{2D^2}\left[4R^2u^2+2u^2(R^2+A)\right]\\
&=\frac{Ru^2}{2D^2}(6R^2+2A)\\
&=\frac{Ru^2}{D},
\end{aligned}\]

与电源端计算完全一致。

\newpage

\section{第 6 题}\label{ux7b2c-6-ux9898}

四块宽度均为 \(l\)
的长黑体板围成正方形长管，外侧绝热。各板均为良导体，板间互不接触，只通过热辐射交换能量。第
1、2、3 块板分别与温度为 \(T_1\)、\(T_2\)、\(T_3\)
的热源接触并保持恒温；各板内表面均为理想黑体，可视为余弦辐射体，即沿与法线夹角为
\(\theta\) 的方向，单位时间、单位面积、单位立体角内的辐射能量正比于
\(\cos\theta\)。

\begin{enumerate}
\def\labelenumi{\arabic{enumi}.}
\tightlist
\item
  对第 4 块板上距左侧边界为 \(x\) 的一点，求其向第 1、2、3
  块板辐射的能量比 \(q_1:q_2:q_3\)。
\item
  求第 4 块板单位长度向第 1、2、3 块板辐射的总能量比 \(Q_1:Q_2:Q_3\)。
\item
  若第 4 块板不与热源接触，求热平衡温度 \(T_4\)。
\end{enumerate}

\begin{figure}[H]
\centering
\includegraphics[width=0.671\linewidth,height=\textheight,keepaspectratio,alt={图 6-1　正方形黑体长管截面及第 4 板上辐射点 P 的几何位置。}]{../figures/question-06.pdf}
\caption*{图 6-1　正方形黑体长管截面及第 4 板上辐射点 \(P\) 的几何位置。}
\end{figure}

\subsection{一、题型定位}\label{ux4e00ux9898ux578bux5b9aux4f4d-5}

三块板的辐射交换要先求局部几何权重，再沿板积分成视因数。热平衡使用温度四次方，投影距离和两面元的实际距离须分别标记。

\subsection{二、知识点}\label{ux4e8cux77e5ux8bc6ux70b9-5}

\begin{itemize}
\tightlist
\item
  黑体余弦辐射定律。
\item
  二维长板结构的视角积分与形系数。
\item
  互易性、归一化与斯忒藩---玻尔兹曼定律。
\end{itemize}

\subsection{三、物理图像}\label{ux4e09ux7269ux7406ux56feux50cf-5}

从第 4
板上一点向内半空间发出的光线最终落到其余三板之一。位置越靠近左侧，第 1
板所张视野越大；位于中点时左右对称。整板视因数就是把每一点的这种分配沿板累加。

\subsection{四、解题入口}\label{ux56dbux89e3ux9898ux5165ux53e3-5}

先选第 4 板上一点
\(P\)，对目标板上的面积元写出源面余弦、接收面余弦和距离平方反比，再按轴向和横向依次积分。

先求第 4
板上一点向各板发出的辐射，再沿整板累加，以吸收功率等于辐射功率求平衡温度。积分结果应满足：中点处左右权重相等，局部权重之和不随位置改变，整板视因数之和为
\(1\)。这些条件可逐步检验下面的计算。

\[\text{点到板的几何权重}\to\text{沿板积分}\to\text{辐射热平衡}.\]

这样积分只是计算权重的工具，不会遮住真正的物理图景。

\subsection{五、解析}\label{ux4e94ux89e3ux6790-5}

图 6-2 左侧标出第 4 板辐射点 \(P\) 到三块接收板面积元 \(dA_1,dA_2,dA_3\)
的三维距离与截面坐标；右侧把纸面外轴向差 \(z\) 明确纳入
\(R^2=\rho^2+z^2\)。第（1）问依次计算三个几何权重并检验和为
\(1\)，第（2）、（3）问再由这些权重列辐射平衡。

\begin{figure}[H]
\centering
\includegraphics[width=0.839\linewidth,height=\textheight,keepaspectratio,alt={图 6-2　第 4 板辐射点到各接收面元的三维投影几何。}]{../figures/solution-06.pdf}
\caption*{图 6-2　第 4 板辐射点到各接收面元的三维投影几何。}
\end{figure}

\textbf{（1）}

按图
6-2，辐射传递权重同时包含源面余弦、接收面余弦和三维距离平方反比。长管虽然用二维截面表示，但面积元还沿轴向延伸，故必须保留
\(z\) 并先作轴向积分。

设余弦辐射强度中的公共比例常数为 \(k\)。取长管轴向为 \(z\) 轴，第 4
板上的辐射点为 \(P\)。在横截面内，\(P\) 到第 1 板的垂直距离为 \(x\)；第
1 板上的横向坐标记为 \(y\)，轴向坐标差记为 \(z\)。因此第 1 板面积元为
\(\mathrm dA_1=\mathrm dy\,\mathrm dz\)，从 \(P\) 到该面元的三维距离满足

\[R^2=x^2+y^2+z^2.\]

源面法线与接收面法线对应的两个方向余弦分别为

\[\cos\theta_s=\frac{y}{R},\qquad
\cos\theta_r=\frac{x}{R},\]

因此

\[\mathrm dq_1
=k\frac{\cos\theta_s\cos\theta_r}{R^2}\,\mathrm dy\,\mathrm dz
=k\frac{xy}{(x^2+y^2+z^2)^2}\,\mathrm dy\,\mathrm dz.\]

先用积分

\[\int_{-\infty}^{+\infty}\frac{\mathrm dz}{(a^2+z^2)^2}
=\frac{\pi}{2a^3},\]

其中 \(a^2=x^2+y^2\)。这个标准积分可由

\[\int_{-\infty}^{+\infty}\frac{\mathrm dz}{a^2+z^2}=\frac\pi a\]

对参数 \(a\) 求导得到：

\[-2a\int_{-\infty}^{+\infty}
\frac{\mathrm dz}{(a^2+z^2)^2}
=-\frac\pi{a^2},\]

故积分值确为 \(\pi/(2a^3)\)。对 \(y\in[0,l]\) 积分时，用到

\[\frac{\mathrm d}{\mathrm dy}
\frac1{\sqrt{x^2+y^2}}
=-\frac{y}{(x^2+y^2)^{3/2}}.\]

于是

\[\begin{aligned}
q_1
&=\frac{\pi kx}{2}
\int_0^l\frac{y\,\mathrm dy}{(x^2+y^2)^{3/2}}\\
&=\frac{\pi kx}{2}
\left[-\frac1{\sqrt{x^2+y^2}}\right]_{0}^{l}\\
&=\frac{\pi k}{2}
\left(1-\frac{x}{\sqrt{x^2+l^2}}\right).
\end{aligned}\]

第 3 板与第 1 板关于中线对称，只需在 \(q_1\) 中作代换 \(x\to l-x\)，得到

\[q_3=\frac{\pi k}{2}
\left(1-\frac{l-x}{\sqrt{(l-x)^2+l^2}}\right).\]

对第 2 板，设横向坐标为 \(\xi\)。点 \(P\) 与该面积元的横截面水平坐标差为
\(\xi-x\)，法向间隔为 \(l\)，轴向坐标差仍为 \(z\)，所以

\[R^2=(\xi-x)^2+l^2+z^2.\]

源面与接收面平行，两个方向余弦均为 \(l/R\)：

\[\mathrm dq_2
=k\frac{l^2}{[(\xi-x)^2+l^2+z^2]^2}\,\mathrm d\xi\,\mathrm dz.\]

依次对 \(z\) 与 \(\xi\) 积分：

\[\begin{aligned}
q_2
&=\frac{\pi kl^2}{2}
\int_0^l
\frac{\mathrm d\xi}{[(\xi-x)^2+l^2]^{3/2}}\\
&=\frac{\pi kl^2}{2}
\left[
\frac{\xi-x}{l^2\sqrt{(\xi-x)^2+l^2}}
\right]_{0}^{l}\\
&=\frac{\pi k}{2}
\left[
\frac{\xi-x}{\sqrt{(\xi-x)^2+l^2}}
\right]_{\xi=0}^{\xi=l}\\
&=\frac{\pi k}{2}\left[
\frac{x}{\sqrt{x^2+l^2}}
+\frac{l-x}{\sqrt{(l-x)^2+l^2}}
\right].
\end{aligned}\]

现在三个局部辐射权重都已由积分得到。去掉公共因子 \(\pi k/2\)，可写成

\[\boxed{
q_1:q_2:q_3=
\left(1-\frac{x}{\sqrt{x^2+l^2}}\right):
\left[
\frac{x}{\sqrt{x^2+l^2}}
+\frac{l-x}{\sqrt{(l-x)^2+l^2}}
\right]:
\left(1-\frac{l-x}{\sqrt{(l-x)^2+l^2}}\right)
}.\]

三者相加恰为

\[q_1+q_2+q_3=\pi k,\]

与 \(x\) 无关，体现了点 \(P\)
向整个内半空间发出的总能量守恒。这是检验积分式的关键条件。

\textbf{（2）}

沿第 4 板对 \(x\in(0,l)\) 积分：

\[Q_i=\int_0^lq_i(x)\,\mathrm dx.\]

利用左右对称性有 \(Q_1=Q_3\)。积分结果为

\[\begin{aligned}
Q_1
&=\frac{\pi k}{2}
\int_0^l\left(1-\frac{x}{\sqrt{x^2+l^2}}\right)\mathrm dx\\
&=\frac{\pi k}{2}
\left[x-\sqrt{x^2+l^2}\right]_{0}^{l}\\
&=\frac{\pi kl}{2}(2-\sqrt2).
\end{aligned}\]

由第（1）问的局部守恒式 \(q_1+q_2+q_3=\pi k\)，沿整块第 4 板积分可得

\[Q_1+Q_2+Q_3
=\int_0^l\pi k\,\mathrm dx
=\pi kl.\]

再利用 \(Q_3=Q_1\)：

\[\begin{aligned}
Q_2
&=\pi kl-Q_1-Q_3\\
&=\pi kl-2\cdot\frac{\pi kl}{2}(2-\sqrt2)\\
&=\pi kl\bigl[1-(2-\sqrt2)\bigr]\\
&=\pi kl(\sqrt2-1).
\end{aligned}\]

因此

\[\boxed{
Q_1:Q_2:Q_3
=(2-\sqrt2):2(\sqrt2-1):(2-\sqrt2)
}.\]

相应的归一化视因数为

\[F_{41}=F_{43}=1-\frac1{\sqrt2},\qquad
F_{42}=\sqrt2-1,\]

并满足 \(F_{41}+F_{42}+F_{43}=1\)。

\textbf{（3）}

三个视因数分别表示第 4 板发出的总辐射中落到第 1、2、3
板的比例。几何视因数满足互易关系

\[A_4F_{4i}=A_iF_{i4}.\]

四块板的面积相等，所以

\[F_{i4}=F_{4i}.\]

这意味着第 \(i\) 板发出的辐射中落到第 4 板的比例，也等于上面求出的
\(F_{4i}\)。单位面积理想黑体在温度 \(T\) 时的总辐射出射度为
\(\sigma T^4\)，故第 4 板从其余三板吸收的功率密度为

\[P_{\mathrm{in}}
=\sigma\left(F_{41}T_1^4+F_{42}T_2^4+F_{43}T_3^4\right).\]

第 4 板自身向整个内腔发出的功率密度为

\[\begin{aligned}
P_{\mathrm{out}}
&=\sigma T_4^4(F_{41}+F_{42}+F_{43})\\
&=\sigma T_4^4.
\end{aligned}\]

热平衡时第 4 板的净辐射功率为零：

\[\sigma T_4^4
=\sigma\left(F_{41}T_1^4+F_{42}T_2^4+F_{43}T_3^4\right).\]

故

\[\boxed{
T_4=
\left[
\left(1-\frac1{\sqrt2}\right)T_1^4
+(\sqrt2-1)T_2^4
+\left(1-\frac1{\sqrt2}\right)T_3^4
\right]^{1/4}
}.\]

亦可保留未归一化权重，写成

\[\boxed{
T_4=
\left[
\frac{(2-\sqrt2)T_1^4
+2(\sqrt2-1)T_2^4
+(2-\sqrt2)T_3^4}{2}
\right]^{1/4}
}.\]

\emph{【拓展视角：交叉弦法】}
对无限长二维封闭腔，整块板之间的视因数还可由 Hottel 交叉弦法直接求得。第
4 板到第 2 板的视因数为
\(\sqrt2-1\)，无需先求点视因数；但本题第（1）问要求位置 \(x\)
的局部比例，因此积分法仍是主线。

图 6-2 左图中的 \(\rho_i\) 是横截面内距离，右图的
\(R=\sqrt{\rho_i^2+z^2}\)
才是源面元与接收面元的三维距离。接收余弦应取接收面的内法线与``接收点指向源点''的方向夹角；若沿入射光的传播方向计算，须补一个负号。正文各核函数用的是正的投影因子。

\newpage

\section{第 7 题}\label{ux7b2c-7-ux9898}

透明样品不同位置的折射率可能不同。Zernike
相衬显微技术通过在傅里叶面处理不同空间频率分量，将样品的相位变化转化为像面的强度变化。

\begin{enumerate}
\def\labelenumi{\arabic{enumi}.}
\tightlist
\item
  一块足够大的平板的振幅透射率为
  \[t(x)=\frac12\left[1+\cos(k_*x)\right].\]
  一束平行光垂直入射，经透镜后在后焦面形成哪几个衍射条纹？
\item
  对透明相位样品，振幅透射率为 \(1\)，但附加小相位
  \[\varphi(x)=\varphi_0t(x),\qquad \varphi_0=10^\circ.\]
  证明相位信息也能反映在衍射后焦面上。
\item
  在傅里叶面将零级斑振幅衰减为原来的 \(1/4\)，并使高频斑的相位增加
  \(\pi/2\)。经过图示 \(4f\) 光路后，求最终像的对比度。
\end{enumerate}

\begin{figure}[H]
\centering
\includegraphics[width=0.9\linewidth,height=\textheight,keepaspectratio,alt={图 7-1　光路结构：样品经第一透镜在傅里叶面分频，经相位板处理后由第二透镜成像。}]{../figures/question-07.pdf}
\caption*{图
7-1　光路结构：样品经第一透镜在傅里叶面分频，经相位板处理后由第二透镜成像。}
\end{figure}

\subsection{一、题型定位}\label{ux4e00ux9898ux578bux5b9aux4f4d-6}

将透明样品的弱相位调制展开成离散空间频谱，在傅里叶面改变零级光相位，最后在像面叠加振幅、取模平方求强度对比度。

\subsection{二、知识点}\label{ux4e8cux77e5ux8bc6ux70b9-6}

\begin{itemize}
\tightlist
\item
  透镜后焦面的傅里叶变换性质。
\item
  余弦光栅的离散空间频谱。
\item
  小相位近似 \(e^{i\varphi}\approx1+i\varphi\)。
\item
  Zernike 相衬中相位调制向振幅调制的转换。
\end{itemize}

\subsection{三、物理图像}\label{ux4e09ux7269ux7406ux56feux50cf-6}

第一透镜把样品分解成零级参考光和 \(\pm k_*\)
侧带。相位板分别处理参考光与侧带，第二透镜再把它们叠加；亮暗来自复振幅干涉，而不是单个频谱斑自身的强弱。

\subsection{四、解题入口}\label{ux56dbux89e3ux9898ux5165ux53e3-6}

第一步不是做傅里叶积分，而是把余弦写成两个复指数。第（2）问只保留
\(e^{i\varphi}\) 的一阶项，使零级和两个侧带直接显现。

先区分三个概念：复振幅是``带相位的光场''，傅里叶面把不同空间频率分开，探测器最后测到的是强度而不是振幅。余弦样品只含
\(0\) 与 \(\pm k_*\) 三个频率，因此是学习相衬的最简单模型。全题路线为

\[\text{余弦分频}\to\text{小相位产生侧带}\to
\text{傅里叶面移相}\to\text{像面干涉产生亮暗}.\]

初学者若能说清``信息已经在侧带中，但还没有变成强度差''，就抓住了相衬显微的核心。

\subsection{五、解析}\label{ux4e94ux89e3ux6790-6}

图 7-2 把傅里叶面中的零级斑和 \(\pm k_*\)
两个侧带分开：零级振幅衰减为原来的 \(1/4\)，两侧带统一增加相位
\(\pi/2\)。第（1）问先求三条谱线的位置与振幅，第（2）问说明未加相位板时强度一阶项消失，第（3）问按图中处理逐项反变换并求像面对比度。

\begin{figure}[H]
\centering
\includegraphics[width=0.9\linewidth,height=\textheight,keepaspectratio,alt={图 7-2　零级斑衰减及 \textbackslash pm k\_* 侧带统一移相。}]{../figures/solution-07.pdf}
\caption*{图 7-2　零级斑衰减及 \(\pm k_*\) 侧带统一移相。}
\end{figure}

\textbf{（1）}

按图 7-2，样品的常数透射项产生 \(k_x=0\) 零级斑，余弦调制分裂成等振幅的
\(\pm k_*\) 两个侧带。三者在傅里叶面分离，才允许相位板分别处理。

由欧拉公式

\[\cos(k_*x)=\frac12\left(e^{ik_*x}+e^{-ik_*x}\right),\]

原透射率逐步改写为

\[\begin{aligned}
t(x)
&=\frac12\left[1+\cos(k_*x)\right]\\
&=\frac12+\frac12\cdot\frac12
\left(e^{ik_*x}+e^{-ik_*x}\right)\\
&=\frac12+\frac14e^{ik_*x}+\frac14e^{-ik_*x}.
\end{aligned}\]

因此后焦面出现三个离散条纹，横向空间频率分别为

\[k_x=0,\qquad k_x=\pm k_*.\]

若入射振幅为 \(A\)，三者的复振幅依次为

\[\frac A2,\qquad \frac A4,\qquad \frac A4.\]

对应衍射角满足

\[\sin\theta_{\pm}=\pm\frac{k_*\lambda}{2\pi}.\]

若透镜焦距为 \(f\)，在近轴条件下，后焦面三个条纹的位置为

\[X_0=0,\qquad
X_{\pm}\approx f\theta_{\pm}
=\pm\frac{f\lambda k_*}{2\pi}.\]

由于光强与复振幅模平方成正比，理想无限光栅的三条谱线强度比为

\[I_0:I_{+1}:I_{-1}
=\left|\frac A2\right|^2:
\left|\frac A4\right|^2:
\left|\frac A4\right|^2
=4:1:1.\]

这里``足够大''的作用是使各级谱线足够窄；若样品只有有限宽度，每个
\(\delta\) 型谱线会被样品孔径的 \(\operatorname{sinc}\) 包络展宽。

\emph{【数学技巧：把余弦写成两项复指数】} 透镜把空间复指数 \(e^{ik_xx}\)
映射到后焦面上的一个确定位置，因此展开式的每一项可直接对应一个衍射级次。

\textbf{（2）}

样品后的复振幅为

\[E_{\mathrm{out}}(x)=Ae^{i\varphi_0t(x)}.\]

由于 \(\varphi_0=10^\circ\approx0.1745\ll1\)，作一阶展开：

\[E_{\mathrm{out}}(x)\approx A\left[1+i\varphi_0t(x)\right].\]

代入 \(t(x)\)：

\[E_{\mathrm{out}}(x)\approx
A\left(1+\frac{i\varphi_0}{2}\right)
+\frac{iA\varphi_0}{4}e^{ik_*x}
+\frac{iA\varphi_0}{4}e^{-ik_*x}.\]

可见后焦面除零级分量外，在 \(\pm k_*\) 处仍出现侧带，其振幅正比于
\(\varphi_0\)，并比未调制参考光多出约 \(\pi/2\)
的相位。因此样品相位信息已被编码到傅里叶面。

更一般地，若透明样品的相位分布为任意小函数 \(\varphi(x)\)，则

\[Ae^{i\varphi(x)}\approx A+iA\varphi(x).\]

傅里叶变换后

\[\widetilde E(k_x)
\approx2\pi A\delta(k_x)
+iA\widetilde\varphi(k_x).\]

第一项是未偏折的零级参考光，第二项就是样品相位分布的空间频谱。这说明结论并不依赖余弦样品；余弦模型只是把一般连续频谱简化为三个离散斑点。

小相位近似的误差为二阶量。由于

\[e^{i\varphi}=1+i\varphi-\frac12\varphi^2+\cdots,\]

当 \(\varphi_0=10^\circ\approx0.1745\) 时，二阶项的最大量级约为

\[\frac{\varphi_0^2}{2}\approx1.5\%,\]

零级被衰减后，相对误差会放大。因此下面区分弱相位近似结果与完整相位模型的数值。

\textbf{（3）}

相位板把零级参考光振幅衰减为 \(A/4\)，并将高频侧带再移相
\(\pi/2\)，使原本与参考光正交的相位扰动转化为同相或反相的振幅扰动。

下面从傅里叶面三个分量逐项还原。保留到 \(\varphi_0\)
的一阶量，第（2）问得到

\[E_{\mathrm{out}}(x)
\approx A\left(1+\frac{i\varphi_0}{2}\right)
+\frac{iA\varphi_0}{4}e^{ik_*x}
+\frac{iA\varphi_0}{4}e^{-ik_*x}.\]

第一项是零级分量。相位板把整个零级分量乘以 \(1/4\)，严格写来是

\[\frac A4\left(1+\frac{i\varphi_0}{2}\right).\]

括号中的 \(i\varphi_0/2\)
只给参考光增加一阶小相位。计算强度的一阶空间调制时，它与不含 \(x\)
的整体相位等价，不改变亮暗对比，因此后面可把参考光写成实振幅 \(A/4\)。

两个侧带原来的系数都是 \(iA\varphi_0/4\)。相位增加 \(\pi/2\)
等价于复振幅再乘

\[e^{i\pi/2}=i.\]

因此每个侧带处理后成为

\[\frac{iA\varphi_0}{4}\times i
=-\frac{A\varphi_0}{4}.\]

第二个透镜作逆傅里叶变换，两侧带在像面相加：

\[\begin{aligned}
-\frac{A\varphi_0}{4}e^{ik_*x}
-\frac{A\varphi_0}{4}e^{-ik_*x}
&=-\frac{A\varphi_0}{4}
\left(e^{ik_*x}+e^{-ik_*x}\right)\\
&=-\frac{A\varphi_0}{2}\cos(k_*x).
\end{aligned}\]

因此原来的相位调制项变成了实振幅调制。忽略不影响一阶对比度的整体相位后，像面复振幅为

\[E_{\mathrm{img}}(x)
\approx A\left[\frac14-\frac{\varphi_0}{2}\cos(k_*x)\right].\]

余弦取 \(\pm1\)
时分别给出两个振幅极值；整体负号只会交换亮暗位置，不改变对比度。先把角度换成弧度：

\[\varphi_0=10^\circ=\frac\pi{18}\approx0.17453,\]

所以像面振幅调制的大小为

\[\frac{\varphi_0}{2}A\approx0.08727A.\]

因此最亮与最暗处的振幅分别为

\[A_{\max}=A(0.25+0.08727),\qquad
A_{\min}=A(0.25-0.08727).\]

像强度与振幅平方成正比。按

\[C=\frac{I_{\max}-I_{\min}}{I_{\max}+I_{\min}}\]

定义对比度，得到

\[C=
\frac{(0.33727)^2-(0.16273)^2}
{(0.33727)^2+(0.16273)^2}
\approx\boxed{0.622}.\]

若保留符号计算，令

\[a=\frac14,\qquad b=\frac{\varphi_0}{2},\]

则

\[I_{\max}=A^2(a+b)^2,\qquad
I_{\min}=A^2(a-b)^2,\]

所以

\[\begin{aligned}
C
&=\frac{(a+b)^2-(a-b)^2}{(a+b)^2+(a-b)^2}\\
&=\frac{4ab}{2(a^2+b^2)}\\
&=\frac{2ab}{a^2+b^2}\\
&=\frac{\varphi_0/4}{1/16+\varphi_0^2/4}
\approx0.622.
\end{aligned}\]

\emph{【拓展视角：参考光的作用】}
相衬显微镜相当于一种空间外差测量：强零级光充当本振，弱相位侧带与它干涉，使原本二阶可见的相位信息转化为一阶强度变化。

\textbf{对比度的精度复核。} 上式保留一阶场再平方，得到
\(C\approx0.622\)；它含部分二阶项，并非严格的一阶强度展开。完整透明相位样品有

\[t(x)=e^{i\varphi_0/2}e^{ib\cos k_*x},\qquad b=\varphi_0/2.\]

零级系数是 \(e^{i\varphi_0/2}J_0(b)\)。若零级乘 \(a=1/4\)、所有非零级乘
\(i\)，像面场准确写成

\[\frac{E_{\rm img}}A=e^{i\varphi_0/2}\left[aJ_0(b)+i\left(e^{ib\cos k_*x}-J_0(b)\right)\right].\]

令 \(z=b\cos k_*x\)，强度为

\[I/A^2=[aJ_0(b)-\sin z]^2+[\cos z-J_0(b)]^2.\]

对 \(\varphi_0=\pi/18\)，极值在
\(z=\mp b\)：\(I_{\max}/A^2=0.1133570\)、\(I_{\min}/A^2=0.0263672\)，故完整结果
\(C=0.6225828\approx0.623\)。若题目按弱相位一阶场近似作答，可报
\(0.622\) 并注明近似。

\newpage

\section{第 8 题}\label{ux7b2c-8-ux9898}

费米黄金定则可作如下概率模型表述：对任意两个量子态
\(\alpha\)、\(\beta\)，体系单位时间内从 \(\alpha\) 态跃迁到 \(\beta\)
态的概率为 \(v_{\alpha\beta}p_\alpha\)，其中 \(p_\alpha\) 为体系当前处于
\(\alpha\) 态的概率，跃迁系数满足

\[v_{\alpha\beta}>0,\qquad v_{\alpha\beta}=v_{\beta\alpha}.\]

\begin{enumerate}
\def\labelenumi{\arabic{enumi}.}
\tightlist
\item
  定义 \[H=\sum_i p_i\ln p_i.\] 求 \(\mathrm dp_i/\mathrm dt\) 与
  \(\mathrm dH/\mathrm dt\)，并证明 \(\mathrm dH/\mathrm dt\leq0\)。
\item
  第 \(i\) 个状态的能量为
  \(E_i\)。设状态数足够多，并要求平均能量守恒，求体系达到稳定状态时各态的概率；结果中允许含待定常数。
\end{enumerate}

\subsection{一、题型定位}\label{ux4e00ux9898ux578bux5b9aux4f4d-7}

对称跃迁的主方程给出 \(H\)
单调不增及其稳态。固定平均能量下的最大熵分布另需约束；二者的假设要分别说明。

\subsection{二、知识点}\label{ux4e8cux77e5ux8bc6ux70b9-7}

\begin{itemize}
\tightlist
\item
  连续时间马尔可夫过程的主方程。
\item
  对称跃迁率与详细平衡。
\item
  对数函数单调性及 \(H\) 定理。
\item
  带归一化和平均能量约束的变分法。
\end{itemize}

\subsection{三、物理图像}\label{ux4e09ux7269ux7406ux56feux50cf-7}

把每个量子态看成一个概率容器，状态之间有双向概率流。若两态概率不同，较大概率一侧的流出更多，概率差逐渐缩小；\(H\)
用来量化这种趋近平衡的过程。

\subsection{四、解题入口}\label{ux56dbux89e3ux9898ux5165ux53e3-7}

先按``流入减流出''写主方程并验证总概率守恒。求 \(\dot H\)
时，关键操作是把原双重求和与交换 \(i,j\) 后的式子相加取平均。

先从两个状态讲起：概率总是从``较多的一边''净流向``较少的一边''，所以分布会逐渐变平。多状态主方程只是把所有两态交换相加。接下来构造
\(H\)，不是为了神奇地猜公式，而是要把 \(\dot H\) 整理成许多项

\[(p_i-p_j)(\ln p_i-\ln p_j),\]

使每一项都能判断符号。第（2）问则是另一条线：在归一化和平均能量约束下寻找
\(H\)
的最小值。教学顺序应是``两态直觉---多态求和---约束极值''，不要一开始就直接写双重求和。

\subsection{五、解析}\label{ux4e94ux89e3ux6790-7}

图 8-1 以三个状态示意对称跃迁系数
\(v_{ij}=v_{ji}\)：两个方向的系数相同，但实际概率流分别为 \(v_{ij}p_i\)
与
\(v_{ij}p_j\)，只有当概率相等时才互相抵消。第（1）问把所有边的净概率流配对求和，证明熵型量
\(H=\sum_i p_i\ln p_i\)
单调不增；第（2）问在归一化和平均能量约束下求其极小分布。

\begin{figure}[H]
\centering
\includegraphics[width=0.9\linewidth,height=\textheight,keepaspectratio,alt={图 8-1　对称跃迁系数下由概率差决定的净概率流。}]{../figures/solution-08.pdf}
\caption*{图 8-1　对称跃迁系数下由概率差决定的净概率流。}
\end{figure}

\textbf{（1）}

按图 8-1，对任意一条连接 \(i,j\)，从 \(i\) 到 \(j\) 的流量与从 \(j\) 到
\(i\) 的流量分别含
\(p_i,p_j\)。跃迁系数对称不等于两股概率流相等；净流总由概率较大的状态指向概率较小的状态。

状态 \(i\) 的概率变化率等于所有流入概率流减去所有流出概率流：

\[\boxed{
\frac{\mathrm dp_i}{\mathrm dt}
=\sum_j\left(v_{ji}p_j-v_{ij}p_i\right)
=\sum_jv_{ij}(p_j-p_i)
}. \tag{1}\]

其中 \(j=i\) 的项恒为零，是否写入求和不影响结果。对任意一对状态
\(i,j\)，从 \(j\) 流向 \(i\) 的概率流为 \(v_{ij}p_j\)，反向概率流为
\(v_{ij}p_i\)，故净流确为 \(v_{ij}(p_j-p_i)\)。

概率归一化守恒，因为

\[\begin{aligned}
\sum_i\frac{\mathrm dp_i}{\mathrm dt}
&=\sum_{i,j}v_{ij}p_j-\sum_{i,j}v_{ij}p_i.
\end{aligned}\]

在第一项中交换哑指标 \(i\leftrightarrow j\)，并使用 \(v_{ji}=v_{ij}\)：

\[\sum_{i,j}v_{ij}p_j
=\sum_{i,j}v_{ji}p_i
=\sum_{i,j}v_{ij}p_i.\]

两项逐项抵消，因此

\[\boxed{\sum_i\frac{\mathrm dp_i}{\mathrm dt}=0}.\]

所以只要初始时 \(\sum_i p_i=1\)，以后任意时刻仍保持归一化。

对 \(H\) 求导：

\[\frac{\mathrm dH}{\mathrm dt}
=\sum_i(1+\ln p_i)\frac{\mathrm dp_i}{\mathrm dt}.\]

由于 \(\sum_i\dot p_i=0\)，

\[\frac{\mathrm dH}{\mathrm dt}
=\sum_i\ln p_i\sum_jv_{ij}(p_j-p_i).\]

交换指标 \(i,j\)，再与原式取平均：

原式为

\[\frac{\mathrm dH}{\mathrm dt}
=\sum_{i,j}v_{ij}(p_j-p_i)\ln p_i.\]

交换哑指标 \(i\leftrightarrow j\)，利用 \(v_{ij}=v_{ji}\)：

\[\frac{\mathrm dH}{\mathrm dt}
=\sum_{i,j}v_{ij}(p_i-p_j)\ln p_j.\]

两式相加后除以 \(2\)：

\[\begin{aligned}
\frac{\mathrm dH}{\mathrm dt}
&=\frac12\sum_{i,j}v_{ij}
\left[(p_j-p_i)\ln p_i+(p_i-p_j)\ln p_j\right]\\
&=\frac12\sum_{i,j}v_{ij}(p_j-p_i)
\left(\ln p_i-\ln p_j\right)\\
&=-\frac12\sum_{i,j}v_{ij}(p_i-p_j)
\ln\frac{p_i}{p_j}.
\end{aligned}\]

\[\boxed{
\frac{\mathrm dH}{\mathrm dt}
=-\frac12\sum_{i,j}v_{ij}(p_i-p_j)
\ln\frac{p_i}{p_j}
}. \tag{2}\]

对任意正数 \(a,b\)，因 \(\ln x\) 单调递增，

\[(a-b)\ln\frac ab\geq0.\]

若不熟悉这一不等式，可以分两种情况直接判断：当 \(a>b\) 时，\(a-b>0\) 且
\(\ln(a/b)>0\)；当 \(a<b\) 时，两因子同时小于零，乘积仍为正；\(a=b\)
时乘积为零。因此每一对状态对应的项都有确定符号。

又因 \(v_{ij}>0\)，故式（2）中每一项均不大于零：

\[\boxed{\frac{\mathrm dH}{\mathrm dt}\leq0}.\]

等号成立当且仅当每一对可连通状态满足
\(p_i=p_j\)，或更一般地满足所有允许跃迁边上的净概率流为零。

若跃迁网络连通，则等号条件迫使所有 \(p_i\)
相等；若网络分成若干互不连通的分量，则只要求每个连通分量内部等概率，而各分量所占总概率由初始条件决定。

对于含零概率的状态，可先取 \(p_i>0\) 推导，再用连续极限 \(p_i\to0^+\)
处理。由于 \(p\ln p\to0\)，\(H\) 本身仍有良好定义。

\textbf{（2）}

稳定分布可由在约束

\[\sum_i p_i=1,\qquad \sum_i p_iE_i=\overline E\]

下极小化 \(H\) 得到。引入拉格朗日乘子 \(\alpha\)、\(\beta\)：

\[\mathcal L=\sum_i p_i\ln p_i
+\alpha\left(\sum_i p_i-1\right)
+\beta\left(\sum_i p_iE_i-\overline E\right).\]

对每个 \(p_i\) 求偏导并令其为零：

\[\frac{\partial\mathcal L}{\partial p_i}
=\ln p_i+1+\alpha+\beta E_i=0.\]

因此

\[p_i=e^{-1-\alpha}e^{-\beta E_i}.\]

令配分函数

\[Z=\sum_i e^{-\beta E_i},\]

把尚未确定的公共因子记为 \(C=e^{-1-\alpha}\)，则

\[p_i=Ce^{-\beta E_i}.\]

代入归一化条件：

\[1=\sum_i p_i
=C\sum_i e^{-\beta E_i}=CZ.\]

所以 \(C=1/Z\)，从而

\[\boxed{p_i=\frac{e^{-\beta E_i}}{Z}}.\]

为确认这是极小值而非极大值，考察二阶变分：

\[\delta^2H=\sum_i\frac{(\delta p_i)^2}{p_i}>0\]

（对任意非零允许变分成立）。因此在归一化与平均能量约束下，所得驻点确为
\(H\) 的唯一极小点，也就是熵 \(S=-kH\) 的唯一极大点。

其中 \(\beta\) 由平均能量条件

\[\overline E=\frac{\sum_iE_ie^{-\beta E_i}}{\sum_i e^{-\beta E_i}}\]

确定。若不存在能量约束且所有状态彼此连通，则
\(\beta=0\)，稳定分布退化为等概率分布。

还可以检查 \(\beta\) 的确定通常是唯一的。记

\[\langle E\rangle_\beta
=\frac{\sum_iE_ie^{-\beta E_i}}{Z},\]

对 \(\beta\) 求导可得

\[\frac{\mathrm d\langle E\rangle_\beta}{\mathrm d\beta}
=-\left(\langle E^2\rangle_\beta-\langle E\rangle_\beta^2\right)
=-\operatorname{Var}_\beta(E)\leq0.\]

因此只要能级并非全部相同，平均能量随 \(\beta\)
单调变化，给定允许范围内的 \(\overline E\) 就能唯一确定 \(\beta\)。

还可写出热力学形式

\[\beta=\frac1{kT},\qquad
p_i=\frac{e^{-E_i/(kT)}}{Z},\]

但本题未给温度，故把 \(\beta\) 保留为由平均能量决定的待定常数更严谨。

需要区分两种模型层次：若题设的对称跃迁允许任意不同能量状态直接互跳，主方程本身只保证等概率稳态；要得到带
\(\beta\neq0\)
的玻尔兹曼分布，还需把``能量守恒''理解为宏观平均能量约束，或把系统与能量守恒的更大封闭体系共同考虑。这里按题目第（2）问的约束变分含义处理。

\textbf{题设两项假设的相容性。} 对任意态均有 \(v_{ij}>0\) 且
\(v_{ij}=v_{ji}\)
时，第（1）问动力学的唯一稳态是等概率分布。它通常不守恒指定的平均能量，因为

\[\frac{\mathrm d\overline E}{\mathrm dt}=\frac12\sum_{ij}v_{ij}(p_j-p_i)(E_i-E_j).\]

例如两态能量为 \(0,\Delta\)，初始 \(p=(1,0)\)，便有
\(\dot{\overline E}=v\Delta>0\)。因此第（2）问的指数分布应作为``给定平均能量的熵最大化问题''解答，不能声称它由原有全连通对称跃迁自动产生。若坚持两项同时成立，只能有特殊能级、特殊初始平均值，或修改允许的跃迁及其平衡条件。

\end{document}
